\documentclass[12pt, a4paper]{article}

\usepackage[utf8]{inputenc}
\usepackage[italian]{babel}
\usepackage{geometry}
\geometry{a4paper, top=3cm, bottom=3cm, left=2.5cm, right=2.5cm}
\usepackage{tabularx}
\usepackage{xltabular}
\usepackage{booktabs}
\usepackage{hyperref}
\usepackage{eurosym}
\usepackage{graphicx}
\usepackage{titlesec}
\usepackage[figure,table]{totalcount}
\usepackage{fancyhdr}
\usepackage{adjustbox}
\usepackage{extramarks}
\usepackage{makecell}
\usepackage{multirow}
\usepackage{array}
\usepackage{float}
\usepackage[skip=10pt plus 2pt, indent=5pt]{parskip}
\usepackage{hyperref}
\usepackage{textgreek}
\usepackage{changepage}
\newcommand{\urlglossario}{https://skynetunigroup.github.io/Code_Guardian/Documentazione/RTB/glossario_v1.0.0.pdf} % Da aggiornare con nuove versioni
\newcommand{\glo}[2]{\href{\urlglossario\#nameddest=#2}{#1\textsubscript{G}}}

\pagestyle{fancy}
\fancyhf{}
\setlength{\headheight}{35pt}

\hypersetup{
    pdftitle={Piano di Progetto},
    colorlinks=true,
    linkcolor=black,
    filecolor=magenta,      
    urlcolor=blue,
}

\newcommand{\doctitle}{Piano di Progetto} % Sostituisci con il titolo reale

% --- INTESTAZIONE (HEADER) ---
\lhead{\nouppercase{\firstleftmark}}
\rhead{\adjustbox{valign=m}{\includegraphics[height=1cm]{logo_skynet_dark.png}}}

% --- PIÈ DI PAGINA (FOOTER) ---
\lfoot{\doctitle} 
\rfoot{\thepage} 

% Linee di separazione
\renewcommand{\headrulewidth}{0.4pt} % Linea sotto l'intestazione
\renewcommand{\footrulewidth}{0.4pt} % Linea sopra il piè di pagina

\begin{document}
% ==========================================
% FRONTESPIZIO
% ==========================================
\begin{titlepage}
    \centering
    \includegraphics[width=0.5\linewidth]{logo_skynet_dark.png}
    
    {\Large Gruppo SkyNet}\\
   
    {\normalsize \href{mailto:skynetunigroup@gmail.com}{skynetunigroup@gmail.com}}
    \vspace*{1.5cm}
    
    {\Huge \textbf{\doctitle}} \\
    
    \vspace{3cm}
    
    % Tabella Informazioni Documento, sostituire con il contenuto reale
    \begin{table}[H]
        \centering
        \renewcommand{\arraystretch}{1.5}
        \begin{tabularx}{\textwidth}{|l|X|}
            \hline
            \textbf{Versione} & \textit{1.4.0} \\ \hline
    
            \textbf{Uso} & \textit{Esterno} \\ \hline
            \textbf{Destinatari} & \textit{Prof. Tullio Vardanega, Prof. Riccardo Cardin} \\ \hline
            \textbf{Redattori} & \textit{Dario Pirrone, Sara Ristovic, Edoardo Granziero, Marco Barbiero, Samuele Bertin, Andrea Fistarol} \\ \hline
            \textbf{Verifica} & \textit{Marco Barbiero, Samuele Bertin, Edoardo Granziero, Dario Pirrone, Sara Ristovic, Andrea Fistarol} \\ \hline
            \textbf{Approvazione} & \textit{Samuele Bertin, Dario Pirrone, Andrea Fistarol} \\ \hline

        \end{tabularx}
    \end{table}
    
    \vfill
\end{titlepage}
\stepcounter{page}

\newpage

% ==========================================
% REGISTRO DELLE VERSIONI
% ==========================================
\section*{Versioni del documento}
\markboth{Versioni del documento}{}
%%%%%%%%%%%%%%%%%
{\renewcommand{\arraystretch}{1.3}
\begin{xltabular}{\textwidth}{|c|c|X|p{3.5cm}|p{3.5cm}|}
    \hline
    \textbf{Ver.} & \textbf{Data} & \textbf{Descrizione} & \textbf{Autore} & \textbf{Verificatore} \\ \hline
    \endfirsthead

    \hline
    \textbf{Ver.} & \textbf{Data} & \textbf{Descrizione} & \textbf{Autore} & \textbf{Verificatore} \\ \hline
    \endhead
        1.4.0 & 2026/09/06 & \textit{Consuntivo} \textit{sprint} 14 e \textit{preventivo} \textit{sprint} 15 & Dario Pirrone & - \\ \hline
        1.3.0 & 2026/08/27 & \textit{Consuntivo} \textit{sprint} 13 e \textit{preventivo} \textit{sprint} 14 & Andrea Fistarol & Edoardo Granziero \\ \hline
        1.2.1 & 2026/08/21 & Aggiunta versione al riferimento del glossario e aggiunta data di accesso alle risorse web di riferimento & Marco Barbiero & Dario Pirrone \\ \hline
        1.2.0 & 2026/08/18 & \textit{Consuntivo} \textit{sprint} 12 e \textit{Preventivo sprint} 13 & Andrea Fistarol & Dario Pirrone \\ \hline
        1.1.0 & 2026/08/13 & \textit{Consuntivo} \textbf{RTB} e \textit{Preventivo sprint} 12 & Andrea Fistarol & Marco Barbiero \\ \hline
        1.0.0 & 2026/08/10 & \textit{Consuntivo sprint} 11 e correzioni per revisione RTB & Andrea Fistarol & Edoardo Granziero \\ \hline
        0.14.0 & 2026/08/07 & Correzione e adeguamento della sezione \textit{3 Pianificazione} & Sara Ristovic & Andrea Fistarol \\ \hline 
        0.13.1 & 2026/08/05 & Implementazione delle correzioni della verifica dettagliata & Sara Ristovic & Samuele Bertin \\ \hline 
        0.13.0 & 2026/08/05 & Scrittura \textit{consuntivo} \textit{sprint} 10 e \textit{preventivo} \textit{sprint} 11 & Dario Pirrone & Sara Ristovic \\ \hline
        0.12.0 & 2026/07/25 & Scrittura \textit{consuntivo} \textit{sprint} 9 e \textit{preventivo} \textit{sprint} 10 & Sara Ristovic & Edoardo Granziero \\ \hline
        0.11.0 & 2026/07/17 & Scrittura \textit{consuntivo} \textit{sprint} 8 e \textit{preventivo} \textit{sprint} 9 & Marco Barbiero & Dario Pirrone \\ \hline
        0.10.0 & 2026/07/09 & Scrittura \textit{consuntivo} \textit{sprint} 7 e \textit{preventivo} \textit{sprint} 8 & Sara Ristovic & Dario Pirrone \\ \hline
        0.9.0 & 2026/07/02 & Scrittura \textit{consuntivo} \textit{sprint} 6 e \textit{preventivo} \textit{sprint} 7 & Samuele Bertin & Sara Ristovic \\ \hline
        0.8.0 & 2026/06/18 & Scrittura \textit{consuntivo} \textit{sprint} 5 e \textit{preventivo} \textit{sprint} 6 & Samuele Bertin & Sara Ristovic \\ \hline
        0.7.0 & 2026/05/29 & Scrittura \textit{consuntivo} \textit{sprint} 4 e \textit{preventivo} \textit{sprint} 5; Descritto il nuovo rischio RO3 & Sara Ristovic & Edoardo Granziero \\ \hline
        0.6.0 & 2026/05/21 & Scrittura capitolo 3: Pianificazione & Edoardo Granziero & Samuele Bertin \\ \hline
        0.5.0 & 2026/05/19 & Scrittura \textit{preventivo} \textit{sprint} 4, scrittura \textit{consuntivo} \textit{sprint} 3 & Edoardo Granziero & Samuele Bertin \\ \hline
        0.4.0 & 2026/05/06 & Scrittura \textit{preventivo} \textit{sprint} 3, scrittura \textit{consuntivo} \textit{sprint} 2, grafico burndown & Samuele Bertin & Dario Pirrone \\ \hline
        0.3.0 & 2026/04/27 & Stesura capitolo 2: Analisi dei Rischi & Dario Pirrone & Marco Barbiero \\ \hline
        0.2.0 & 2026/04/22 & Scrittura \textit{preventivo} \textit{sprint} 1 e \textit{sprint} 2, scrittura \textit{consuntivo} \textit{sprint} 1 & Dario Pirrone & Marco Barbiero \\ \hline
        0.1.0 & 2026/04/12 & Scrittura della sezione ``Introduzione'' & Sara Ristovic & Marco Barbiero \\ \hline
        0.0.1 & 2026/04/11 & Strutturazione del documento & Dario Pirrone & Marco Barbiero \\ \hline

\end{xltabular}
}

\newpage

% ==========================================
% INDICI
% ==========================================

\tableofcontents
\iftotaltables  % Se sono presenti tabelle o figure appare l'indice
    \vspace{1cm}
    \listoftables
\fi

\iftotalfigures
    \vspace{1cm}
    \listoffigures
\fi
\newpage

% ==========================================
% CONTENUTO DEL DOCUMENTO
% ==========================================

\section{Introduzione}

\subsection{Scopo del documento}

Questo documento ha lo scopo di pianificare e monitorare i progressi nello sviluppo del progetto ``\textbf{\glo{Code Guardian}{code-guardian}}''.
In particolare serve a tracciare:

\begin{itemize}
\item
  le attività da svolgere e quelle già completate
\item
  i rischi individuati e la loro gestione
\item
  la pianificazione delle \textit{\glo{milestone}{milestone}} del progetto (a grana grossa) e dei
  singoli \textit{\glo{sprint}{sprint}} (a grana fine)
\item
  i costi temporali ed economici, sia previsti che effettivi.
\end{itemize}

L\textquotesingle adozione di tale documento mira a incrementare
l\textquotesingle efficacia e l\textquotesingle efficienza del gruppo di
lavoro.

Un ulteriore scopo importante di questo documento è di servire come base
per riflettere sull\textquotesingle andamento di ogni \textit{sprint},
analizzando le difficoltà incontrate e lo stato degli obiettivi per
capire come migliorarci nello \textit{sprint} successivo.

Il documento verrà aggiornato continuamente durante tutto il periodo di
lavoro e sarà completato solo alla fine del progetto.

\subsection{\textbf{Glossario}}

Termini tecnici e vocaboli particolarmente importanti il cui
significato non sia immediatamente noto o chiaro vengono definiti nel
documento \textbf{Glossario} che è in costante aggiornamento.
Ogni termine
presente nel \textbf{Glossario} è contrassegnato da una G a pedice ed è un
collegamento ipertestuale che rimanda direttamente alla definizione
corrispondente nel \textbf{Glossario}.
Tali parole vengono contrassegnate
solamente alla loro prima occorrenza in questo documento.

I termini evidenziati in grassetto identificano concetti tecnici ed elementi specifici del progetto, mentre i termini in corsivo indicano concetti tecnici di carattere generale, comuni al dominio dello sviluppo software.

\subsection{Riferimenti}

\subsubsection{Riferimenti normativi}

\begin{itemize}
\item
  \textbf{\glo{Norme di Progetto}{norme-di-progetto}}
\item  
  \href{https://www.math.unipd.it/~tullio/IS-1/2025/Dispense/PD1.pdf}{Regolamento Progetto Didattico} (consultato il 10/08/2026)
\item
  \href{https://www.math.unipd.it/~tullio/IS-1/2025/Progetto/C2.pdf}{Capitolato C2 - \textbf{Code Guardian}} (consultato il 10/08/2026)
\end{itemize}

\subsubsection{Riferimenti informativi}

\begin{itemize}
\item
  \href{https://skynetunigroup.github.io/Code_Guardian/Documentazione/RTB/glossario_v1.0.0.pdf}{\textbf{Glossario} v1.0.0}
\item
  \href{https://www.math.unipd.it/~tullio/IS-1/2025/Dispense/T02.pdf}{Slide T02 del corso Ingegneria del Software (Processi di Ciclo di Vita)} (consultato il 10/08/2026)
\item
    \href{https://www.math.unipd.it/~tullio/IS-1/2025/Dispense/T03.pdf}{Slide T03 del corso Ingegneria del Software (Il Ciclo di Vita del SW)} (consultato il 10/08/2026)
\item
  \href{https://www.math.unipd.it/~tullio/IS-1/2025/Dispense/T04.pdf}{Slide T04 del corso Ingegneria del Software (Gestione di Progetto)} (consultato il 10/08/2026)

\end{itemize}

\newpage
\section{Analisi dei Rischi}

L\textquotesingle attività di gestione dei rischi è un processo
fondamentale e continuo che accompagna l\textquotesingle intero ciclo di
vita del progetto.
L\textquotesingle obiettivo principale è identificare
preventivamente i potenziali ostacoli che potrebbero compromettere il
rispetto delle scadenze, la qualità del prodotto o il budget orario
preventivato.

Il gruppo adotta un approccio proattivo: per ogni rischio individuato
viene definita una strategia di precauzione, volta a minimizzare la probabilità che l'evento si verifichi, e una strategia di mitigazione, per ridurre l'impatto qualora il rischio diventi realtà.

Per garantire una gestione strutturata, i rischi sono stati classificati
nelle seguenti categorie:

\begin{itemize}
\item
  \textbf{RO (Rischi Organizzativi):} legati alla pianificazione, al
  coordinamento e alla logistica del team.
\item
  \textbf{RT (Rischi Tecnologici):} relativi all\textquotesingle uso di
  strumenti, framework e alla complessità tecnica.
\item
  \textbf{RI (Rischi Individuali):} inerenti ai singoli componenti e
  alle loro disponibilità o competenze.
\item
  \textbf{RR (Rischi legati ai Requisiti):} legati alla comprensione e
  alla variazione delle specifiche di progetto.
\end{itemize}

Ogni rischio è valutato secondo due parametri principali:

\begin{enumerate}
\def\labelenumi{\arabic{enumi}.}
\item
  \textbf{Probabilità:} la frequenza attesa con cui il rischio può
  manifestarsi (Bassa, Media, Alta).
\item
  \textbf{Pericolosità:} l\textquotesingle entità del danno che il
  rischio arrecherebbe al progetto (Bassa, Media, Alta).
\end{enumerate}

L\textquotesingle analisi non è da considerarsi statica: il
\textbf{\glo{Responsabile}{responsabile}} di Progetto si occupa di monitorare costantemente
l\textquotesingle insorgenza di nuove criticità e di aggiornare i
\textit{\glo{Consuntivi}{consuntivo}} di periodo in base all\textquotesingle efficacia delle
contromisure adottate.
\newpage

\subsection{Rischi organizzativi}
\begin{table}[H]
    \centering
    \renewcommand{\arraystretch}{1.3}
    \begin{tabularx}{\textwidth}{|c|X|}
        \hline
        \multicolumn{2}{|c|}{\textbf{Scarsa Pianificazione - RO1}} \\ \hline
        Descrizione & Difficoltà nell\textquotesingle individuare correttamente le attività
        necessarie e nel suddividere i compiti in modo efficiente, causando
        potenziali ritardi e spreco di risorse.
\\ \hline
        Probabilità di occorrenza & Alta.
\\ \hline
        Grado di pericolosità & Alta.
\\ \hline
        Precauzioni & Adottare inizialmente una pianificazione pessimistica, definire attività, compiti ed obiettivi molto chiari e con un monitoraggio costante dell'andamento del progetto tramite \textit{Consuntivi} di periodo. Discutere collegialmente le modalità operative prima dell'esecuzione.
\\ \hline
        Mitigazioni possibili & 
        Se la distanza tra \textit{\glo{Preventivo}{preventivo}} e \textit{Consuntivo} di uno \textit{sprint} è eccessiva, il
        team deve ripianificare gli \textit{sprint} successivi prima di procedere con
        altre attività e, se necessario, ridefinire il \textit{\glo{Way of Working}{wow}}.\\ \hline
    \end{tabularx}
    \caption{Scarsa Pianificazione - RO1}
\end{table}


\begin{table}[H]
    \centering
    \renewcommand{\arraystretch}{1.3}
    \begin{tabularx}{\textwidth}{|c|X|}
        \hline
  
        \multicolumn{2}{|c|}{\textbf{Comunicazione Interna Inefficace - RO2}} \\ \hline
        Descrizione & Mancanza di regolarità o chiarezza nelle comunicazioni tra i membri, con
        rischio di fraintendimenti, duplicazione del lavoro o discrepanze nelle
        aspettative.
\\ \hline
        Probabilità di occorrenza & Media-Bassa.
\\ \hline
        Grado di pericolosità & Media.
\\ \hline
        Precauzioni & Definire formalmente canali e tecnologie di comunicazione nelle \textbf{Norme di Progetto} e organizzare incontri interni periodici per confronti diretti.
\\ \hline
        Mitigazioni possibili & Garantire aggiornamenti costanti sullo stato di avanzamento;
incoraggiare l\textquotesingle uso frequente dei canali interni per
        risolvere dubbi non appena si presentano.
\\ \hline
    \end{tabularx}
    \caption{Comunicazione Interna Inefficace - RO2}
\end{table}

\begin{table}[H]
    \centering
    \renewcommand{\arraystretch}{1.3}
    \begin{tabularx}{\textwidth}{|c|X|}
        \hline
        \multicolumn{2}{|c|}{\textbf{Comunicazione Esterna Inefficace/Inefficiente - RO3}} \\ \hline
        Descrizione & 
        Difficoltà nella comunicazione con il \textbf{Proponente}, dovute a risposte lente, messaggi poco chiari o ritardi nell'organizzazione degli incontri, che possono rallentare la risoluzione di dubbi progettuali e le attività dipendenti da riscontri esterni.
\\ \hline
        Probabilità di occorrenza & Alta.
\\ \hline
        Grado di pericolosità & Alta.
\\ \hline
        Precauzioni &
        Formulare richieste al \textbf{Proponente} chiare, complete e tracciabili, preferendo comunicazioni scritte per le decisioni importanti; pianificare gli incontri con sufficiente anticipo.
\\ \hline
        Mitigazioni possibili & 
        Sollecitare il \textbf{Proponente} evidenziando la necessità del riscontro; spostare la priorità su compiti non bloccati dalla comunicazione esterna; se necessario, procedere temporaneamente su assunti ragionevoli, esplicitandoli nella documentazione.
\\ \hline
        
    \end{tabularx}
    \caption{Comunicazione Esterna Inefficace/Inefficiente - RO3}
\end{table}

\begin{table}[H]
    \centering
    \renewcommand{\arraystretch}{1.3}
    \begin{tabularx}{\textwidth}{|c|X|}
        \hline
        \multicolumn{2}{|c|}{\textbf{Sottostima delle Attività - RO4}} \\ \hline
        Descrizione &
        Il tempo effettivamente richiesto per una task si rivela superiore a quello preventivato, portando a ritardi sulle scadenze.
        \\ \hline
        Probabilità di occorrenza & Alta.
        \\ \hline
        Grado di pericolosità & Alta.
        \\ \hline
        Precauzioni &
        Analisi dettagliata dei requisiti e dei compiti prima della stima;
        inclusione di margini di tolleranza nei tempi e pianificazione pessimistica.
        \\ \hline
        Mitigazioni possibili & 
        Segnalazione tempestiva del ritardo da parte del membro assegnatario;
        il team valuta la redistribuzione del lavoro tra più persone o il posticipo di attività meno critiche.\\ \hline
    \end{tabularx}
    \caption{Sottostima delle Attività - RO4}
\end{table}

\begin{table}[H]
    \centering
    \renewcommand{\arraystretch}{1.3}
    \begin{tabularx}{\textwidth}{|c|X|}
        \hline
        \multicolumn{2}{|c|}{\textbf{Sovrastima delle Attività - RO5}} \\ \hline
        Descrizione &
        Il tempo richiesto per completare una 
task è sensibilmente inferiore al
        previsto, rischiando un utilizzo inefficiente delle ore produttive
        rimaste e del tempo rimasto nello \textit{Sprint}.
\\ \hline
        Probabilità di occorrenza & Media.
\\ \hline
        Grado di pericolosità & Bassa.
\\ \hline
        Precauzioni & Valutazione più accurata della complessità delle task basandosi sull'esperienza acquisita negli \textit{sprint} precedenti.
\\ \hline
        Mitigazioni possibili & 
        Comunicazione immediata del completamento anticipato;
il membro libero
        fornisce supporto a chi è in ritardo o avvia nuove task pianificate per
        periodi successivi.\\ \hline
    \end{tabularx}
    \caption{Sovrastima delle Attività - RO5}
\end{table}

\begin{table}[H]
    \centering
    \renewcommand{\arraystretch}{1.3}
    \begin{tabularx}{\textwidth}{|c|X|}
        \hline
        \multicolumn{2}{|c|}{\textbf{Forte Disaccordo nel Gruppo - RO6}} \\ \hline
        Descrizione & 
        Conflitti tra i 
componenti su scelte progettuali, tecnologiche o
        organizzative che bloccano il processo decisionale.
\\ \hline
        Probabilità di occorrenza & Bassa.
\\ \hline
        Grado di pericolosità & Media.
\\ \hline
        Precauzioni &
        Promuovere un clima di collaborazione e ascolto;
definire tempi
        prestabiliti per le discussioni.
\\ \hline
        Mitigazioni possibili & 
        Presentazione formale delle diverse proposte con relativa analisi di pro
        e contro;
valutazione collegiale e, se necessario, ricorso a votazione
        anonima per individuare la scelta più condivisa.\\ \hline
    \end{tabularx}
    \caption{Forte Disaccordo nel Gruppo - RO6}
\end{table}

\begin{table}[H]
    \centering
    \renewcommand{\arraystretch}{1.3}
    \begin{tabularx}{\textwidth}{|c|X|}
        \hline
        \multicolumn{2}{|c|}{\textbf{Problemi di Forza Maggiore - RO7}} \\ \hline
        Descrizione & 
        Eventi esterni imprevedibili, per esempio guasti tecnici o emergenze varie, che rallentano lo sviluppo del progetto e distorcono le previsioni.
\\ \hline
        Probabilità di occorrenza & Media-Alta.
\\ \hline
        Grado di pericolosità & Alta.
\\ \hline
        Precauzioni & 
        Mantenere una pianificazione flessibile e pessimistica che permetta di
        assorbire imprevisti minori.
\\ \hline
        Mitigazioni possibili &
        Monitoraggio costante della situazione;
comunicazione tempestiva a tutti
        i membri del team e possibile ridefinizione delle scadenze interne.
\\ \hline
    \end{tabularx}
    \caption{Problemi di Forza Maggiore - RO7}
\end{table}

\begin{table}[H]
    \centering
    \renewcommand{\arraystretch}{1.3}
    \begin{tabularx}{\textwidth}{|c|X|}
        \hline
        \multicolumn{2}{|c|}{\textbf{Riunioni inefficaci - RO8}} \\ \hline
        Descrizione & 
        Le riunioni falliscono nel raggiungere scopo, allineamento, lasciano questioni in sospeso. Questo può causare mancanza di coordinamento, disallineamento in fini e modalità, ritardare la risoluzione di problemi e costringere all'uso di canali subottimali.
        \\ \hline
        Probabilità di occorrenza & Media-Alta.
        \\ \hline
        Grado di pericolosità & Alta.
        \\ \hline
        Precauzioni & 
        Adottare un Ordine del Giorno (OdG) per la discussione e affrontare a fondo tutti i punti; preparazione individuale rispetto all'OdG; liberarsi da impegni personali per presenziare alla riunione nella sua completezza.
        \\ \hline
        Mitigazioni possibili &
        Programmare una riunione di recupero; utilizzare canali testuali per questioni minori.
        \\ \hline
    \end{tabularx}
    \caption{Riunioni inefficaci - RO8}
\end{table}

%\newpage

\subsection{Rischi tecnologici}

\begin{table}[H]
    \centering
    \renewcommand{\arraystretch}{1.3}
    \begin{tabularx}{\textwidth}{|c|X|}
        \hline
        \multicolumn{2}{|c|}{\textbf{Tecnologie Sconosciute - RT1}} \\ \hline
        Descrizione &
        Difficoltà derivanti dall\textquotesingle utilizzo di tecnologie,
        strumenti di supporto, framework o linguaggi nuovi o particolarmente
        complessi per i membri del gruppo 
(es. \textit{\glo{React}{react}}, \textit{\glo{OWASP}{owasp}}, \textit{\glo{Node.js}{node-js}}). \\ \hline
        Probabilità di occorrenza & Alta.
\\ \hline
        Grado di pericolosità & Alta.
\\ \hline
        Precauzioni & 
        Ogni membro deve dichiarare il proprio livello di conoscenza prima della
        scelta tecnologica.
Prevedere periodi di formazione individuale e
        collettiva prima dell\textquotesingle inizio delle attività di codifica.
\\ \hline
        Mitigazioni possibili & 
        Studio approfondito di materiale esterno;
supporto reciproco tra i
        membri del team;
richiesta di chiarimenti o sessioni di supporto
        all\textquotesingle azienda \textbf{\glo{Proponente}{proponente}}.
\\ \hline
    \end{tabularx}
    \caption{Tecnologie Sconosciute - RT1}
\end{table}

\begin{table}[H]
    \centering
    \renewcommand{\arraystretch}{1.3}
    \begin{tabularx}{\textwidth}{|c|X|}
        \hline
        \multicolumn{2}{|c|}{\textbf{Mancanza di Documentazione - RT2}} \\ \hline
        Descrizione &
        Le tecnologie scelte potrebbero avere una documentazione limitata,
        obsoleta o poco chiara, rendendo difficile la risoluzione di \textbf{\glo{bug}{bug}} o
        l\textquotesingle implementazione di funzionalità 
specifiche. \\ \hline
        Probabilità di occorrenza & Media-Bassa.
\\ \hline
        Grado di pericolosità & Media-Alta.
\\ \hline
        Precauzioni & 
        Verificare l\textquotesingle ampiezza della documentazione e il supporto
        della community prima di adottare definitivamente una tecnologia.
\\ \hline
        Mitigazioni possibili &
        Utilizzare canali di comunicazione interni per condividere scoperte;
cercare guide di terze parti; rivolgersi all\textquotesingle azienda
        \textbf{Proponente} per ottenere materiale aggiuntivo o consigliato.\\ \hline
    \end{tabularx}
    \caption{Mancanza di Documentazione - RT2}
\end{table}

\begin{table}[H]
    \centering
    \renewcommand{\arraystretch}{1.3}
    \begin{tabularx}{\textwidth}{|c|X|}
        \hline
        \multicolumn{2}{|c|}{\textbf{Errori nel Codice - RT3}} \\ \hline
        Descrizione &
        L\textquotesingle attività di codifica e i relativi test possono far
      
        emergere comportamenti inattesi o malfunzionamenti che richiedono
        revisioni e correzioni impreviste.
\\ \hline
        Probabilità di occorrenza & Alta.
\\ \hline
        Grado di pericolosità & Media.
\\ \hline
        Precauzioni & 
        Definizione rigorosa delle \textbf{Norme di Progetto} per la codifica;
adozione
        di \textit{\glo{Test di Unità}{test-di-unita}} e \textit{\glo{Test di Integrazione}{test-di-integrazione}} sin dalle prime fasi di sviluppo.
\\ \hline
        Mitigazioni possibili &
        Il \textbf{\glo{Programmatore}{programmatore}} assegna priorità alla risoluzione del bug;
se la
        criticità è elevata, il \textbf{Responsabile} alloca risorse aggiuntive o
        richiede il supporto dei membri più esperti.\\ \hline
    \end{tabularx}
    \caption{Errori nel Codice - RT3}
\end{table}

\begin{table}[H]
    \centering
    \renewcommand{\arraystretch}{1.3}
    \begin{tabularx}{\textwidth}{|c|X|}
        \hline
        \multicolumn{2}{|c|}{\textbf{Inefficacia degli Strumenti di Supporto - RT4}} \\ \hline
        Descrizione &
        Gli strumenti scelti 
per il coordinamento (editor collaborativi,
        piattaforme di comunicazione, \textit{\glo{tool}{tool}} di gestione \textit{\glo{Repository}{repository}}) si rivelano
        inadatti alle dinamiche del team o presentano limiti tecnici che
        ostacolano la produttività.
\\ \hline
        Probabilità di occorrenza & Media-Bassa.
\\ \hline
        Grado di pericolosità & Bassa.
\\ \hline
        Precauzioni &
        Effettuare una valutazione preliminare dell'efficacia
        degli strumenti attraverso un periodo di prova limitato prima di
        renderli ufficiali nelle \textbf{Norme di Progetto}.
\\ \hline
        Mitigazioni possibili &
        Discussione collegiale sui problemi riscontrati; Se
        l'inefficienza è confermata, l'\textbf{Amministratore} deve
        disporre la sostituzione tempestiva dello strumento con uno più idoneo.\\ \hline
    \end{tabularx}
    \caption{Inefficacia degli Strumenti di Supporto - RT4}
\end{table}

\begin{table}[H]
    \centering
    \renewcommand{\arraystretch}{1.3}
    \begin{tabularx}{\textwidth}{|c|X|}
        \hline
        \multicolumn{2}{|c|}{\textbf{Incompatibilità Tecnologica - RT5}} \\ \hline
        Descrizione &
        Problemi di integrazione tra diverse librerie, framework o servizi
        esterni scelti che non riescono a comunicare correttamente tra loro,
        rallentando o bloccando lo sviluppo dell'architettura.
\\ \hline
        Probabilità di occorrenza & Media.
\\ \hline
        Grado di pericolosità & Media-Alta.
\\ \hline
        Precauzioni &
        Realizzazione di un \textbf{\glo{Proof of Concept (PoC)}{poc}} per testare la fattibilità dell'integrazione tra i componenti critici.
\\ \hline
        Mitigazioni possibili &
        Rivalutazione dell'architettura e sostituzione del componente incompatibile con uno alternativo già testato o più adatto.\\ \hline
    \end{tabularx}
    \caption{Incompatibilità Tecnologica - RT5}
\end{table}

%\newpage

\subsection{Rischi individuali}

\begin{table}[H]
    \centering
    \renewcommand{\arraystretch}{1.3}
    \begin{tabularx}{\textwidth}{|c|X|}
        \hline
        \multicolumn{2}{|c|}{\textbf{Indisponibilità e Vincoli di Tempo - RI1}} \\ \hline
        Descrizione &
 
        Possibili periodi di assenza o riduzione della produttività di uno o più
        membri dovuti a impegni accademici, motivi di salute o imprevisti
        personali.
\\ \hline
        Probabilità di occorrenza & Alta.
\\ \hline
        Grado di pericolosità & Media-Alta.
\\ \hline
        Precauzioni &
        Ogni membro deve comunicare tempestivamente i propri impegni futuri per
        permettere una pianificazione basata sulla disponibilità reale.
\\ \hline
        Mitigazioni possibili &
        Il \textbf{Responsabile} provvede alla redistribuzione delle task del membro
        indisponibile tra gli altri componenti;
in casi critici, si procede allo
        slittamento di attività non prioritarie.\\ \hline
    \end{tabularx}
    \caption{Indisponibilità e Vincoli di Tempo - RI1}
\end{table}

\begin{table}[H]
    \centering
    \renewcommand{\arraystretch}{1.3}
    \begin{tabularx}{\textwidth}{|c|X|}
        \hline
        \multicolumn{2}{|c|}{\textbf{Studio Insufficiente di un Argomento - RI2}} \\ \hline
        Descrizione &
        Un membro affronta lo sviluppo di una funzionalità o la stesura di un documento senza aver acquisito le competenze necessarie, portando a un prodotto di scarsa qualità o errato.
\\ \hline
        Probabilità di occorrenza & Media-Alta.
\\ \hline
        Grado di pericolosità & Media-Alta.
\\ \hline
        Precauzioni &
        Definire dei tempi obbligatori di auto-formazione prima
        dell\textquotesingle inizio di ogni nuova fase tecnologica.
\\ \hline
        Mitigazioni possibili &
        \textit{\glo{Peer Review}{peer-review}} (Revisione incrociata) immediata;
se l\textquotesingle errore
        persiste, il membro viene affiancato da un componente più esperto o si
        richiede un incontro di chiarimento con il \textbf{Proponente}.\\ \hline
    \end{tabularx}
    \caption{Studio Insufficiente di un Argomento - RI2}
\end{table}

\begin{table}[H]
    \centering
    \renewcommand{\arraystretch}{1.3}
    \begin{tabularx}{\textwidth}{|c|X|}
        \hline
        \multicolumn{2}{|c|}{\textbf{Problemi Interpersonali e Conflitti - RI3}} \\ \hline
        Descrizione &
     
        Nascita di tensioni, incomprensioni o antipatie personali tra membri del
        gruppo che possono minare il clima collaborativo e la produttività.
\\ \hline
        Probabilità di occorrenza & Bassa.
\\ \hline
        Grado di pericolosità & Alta.
\\ \hline
        Precauzioni & Promuovere una cultura di feedback costruttivo e trasparenza.
\\ \hline
        Mitigazioni possibili &
        Mediazione diretta del \textbf{Responsabile} per risolvere il conflitto;
se
        necessario, assegnare i membri in conflitto a task o sottogruppi diversi
        per ridurre le frizioni dirette.\\ \hline
    \end{tabularx}
    \caption{Problemi Interpersonali e Conflitti - RI3}
\end{table}

\begin{table}[H]
    \centering
    \renewcommand{\arraystretch}{1.3}
    \begin{tabularx}{\textwidth}{|c|X|}
        \hline
        \multicolumn{2}{|c|}{\textbf{Calo di Motivazione o Impegno - RI4}} \\ \hline
        Descrizione & 
        Disinteresse progressivo 
di un membro verso il progetto, che si
        manifesta con scarsa partecipazione alle riunioni o mancato rispetto
        delle scadenze individuali.
\\ \hline
        Probabilità di occorrenza & Bassa.
\\ \hline
        Grado di pericolosità & Media.
\\ \hline
        Precauzioni &
        Coinvolgimento attivo di tutti i membri nelle scelte decisionali;
valorizzazione del lavoro svolto attraverso gli incontri interni di chiusura dello \textit{Sprint}.
\\ \hline
        Mitigazioni possibili & 
        Colloquio privato con il membro per identificare le cause del
        disinteresse;
se l\textquotesingle impegno non migliora, ridistribuzione
        del carico di lavoro e segnalazione nel verbale interno.\\ \hline
    \end{tabularx}
    \caption{Calo di Motivazione o Impegno - RI4}
\end{table}

%\newpage

\subsection{Rischi legati ai requisiti}

\begin{table}[H]
    \centering
    \renewcommand{\arraystretch}{1.3}
    \begin{tabularx}{\textwidth}{|c|X|}
        \hline
        \multicolumn{2}{|c|}{\textbf{Analisi dei Requisiti Incompleta o Errata - RR1}} \\ \hline
        Descrizione &
        Durante la fase di analisi, il gruppo 
potrebbe omettere dei requisiti
        fondamentali o definirli in modo non corretto, portando a un prodotto
        finale che non soddisfa le attese.
\\ \hline
        Probabilità di occorrenza & Media.
\\ \hline
        Grado di pericolosità & Alta.
\\ \hline
        Precauzioni &
        Studio approfondito del capitolato e dei verbali esterni; utilizzo di \textit{\glo{template}{template}} standard per la definizione dei requisiti e dei \textit{\glo{Casi d\textquotesingle uso}{use-case}}.
\\ \hline
        Mitigazioni possibili & 
        Revisione continua del documento di \textbf{\glo{Analisi dei Requisiti}{analisi-dei-requisiti}};
confronto
        diretto e frequente con il \textbf{Proponente} per validare ogni requisito
        individuato.\\ \hline
    \end{tabularx}
    \caption{\textbf{Analisi dei Requisiti} Incompleta o Errata - RR1}
\end{table}

\begin{table}[H]
    \centering
    \renewcommand{\arraystretch}{1.3}
    \begin{tabularx}{\textwidth}{|c|X|}
        \hline
        \multicolumn{2}{|c|}{\textbf{Cambiamento dei Requisiti in Itinere - RR2}} \\ \hline
        Descrizione &
        Richiesta di modifiche o aggiunte ai 
requisiti da parte del \textbf{Proponente}
        durante lo svolgimento del progetto, o necessità di modifica emersa a
        causa di vincoli tecnologici scoperti in ritardo.
\\ \hline
        Probabilità di occorrenza & Media-Bassa.
\\ \hline
        Grado di pericolosità & Media.
\\ \hline
        Precauzioni &
        Definizione chiara dei requisiti \textit{\glo{must have}{must-have}};
avvisare il \textbf{Proponente}
        dell\textquotesingle impatto che ogni modifica ha sui costi e sui tempi
        di consegna.
\\ \hline
        Mitigazioni possibili &
        Valutazione dell\textquotesingle impatto della modifica sul \textbf{\glo{Piano di Progetto}{piano-di-progetto}};
se la modifica è critica, il \textbf{Responsabile} deve ripianificare
        lo \textit{Sprint} corrente e quelli futuri; se le risorse rimaste non permettono di implementare i cambiamenti richiesti, si procede con la negoziazione delle modifiche con il \textbf{Proponente}.\\ \hline
    \end{tabularx}
    \caption{Cambiamento dei Requisiti in Itinere - RR2}
\end{table}

{\begin{table}[H]
    \centering
    \renewcommand{\arraystretch}{1.3}
    \begin{tabularx}{\textwidth}{|c|X|}
        \hline
        \multicolumn{2}{|c|}{\textbf{Eccessivo Numero di Requisiti Opzionali - RR3}} \\ \hline
        Descrizione & 
        Il gruppo identifica e si impegna a realizzare troppi requisiti
 
        desiderabili o opzionali, rischiando di non terminare in tempo quelli
        obbligatori.
\\ \hline
        Probabilità di occorrenza & Media-Bassa.
\\ \hline
        Grado di pericolosità & Media.
\\ \hline
        Precauzioni & 
        Classificare rigorosamente i requisiti in "Obbligatori", "Desiderabili"
        e "Opzionali".
\\ \hline
        Mitigazioni possibili &
        In caso di ritardo sulla tabella di marcia, il \textbf{Responsabile} deve
        disporre l'abbandono immediato dei requisiti opzionali
        per concentrare tutte le risorse sul nucleo essenziale del software.\\ \hline
    \end{tabularx}
    \caption{Eccessivo Numero di Requisiti Opzionali - RR3}
\end{table}

\begin{table}[H]
    \centering
    \renewcommand{\arraystretch}{1.3}
    \begin{tabularx}{\textwidth}{|c|X|}
        \hline
     
        \multicolumn{2}{|c|}{\textbf{Conflitto tra Requisiti - RR4}} \\ \hline
        Descrizione &
        Presenza di due o più requisiti che si contraddicono a vicenda in
        qualche forma.
\\ \hline
        Probabilità di occorrenza & Bassa.
\\ \hline
        Grado di pericolosità & Bassa.
\\ \hline
        Precauzioni &
        Creazione di una \textit{\glo{matrice di tracciabilità}{matrice-di-tracciabilita}} dei requisiti per individuare
        dipendenze e potenziali conflitti.
\\ \hline
        Mitigazioni possibili &
        Discussione collegiale e consultazione con il \textbf{Proponente} per stabilire
        quale dei requisiti in conflitto abbia la priorità o per trovare un
        compromesso accettabile.
\\ \hline
    \end{tabularx}
    \caption{Conflitto tra Requisiti - RR4}
\end{table}

\newpage
\section{Pianificazione}

La pianificazione delle attività è stata definita per promuovere uno sviluppo strutturato e tempestivo del prodotto, il cui obiettivo principale è la realizzazione di tre agenti: agente \textit{\glo{README}{readme}}, agente Security e agente \textit{\glo{Changelog}{changelog}}.
Il periodo di lavoro ufficiale si estende dal 10 marzo fino al limite finale previsto per il 21 settembre. Per gestire al meglio le risorse e rispettare le scadenze, le attività sono state suddivise in periodi di lavoro ben definiti.

\subsection{\textbf{Macro-fasi} del progetto}
Il progetto è stato diviso in due \textbf{\glo{macro-fasi}{macro-fasi}} principali, ognuna associata a una \textit{milestone} di revisione fondamentale per valutare l'avanzamento dei lavori e la qualità dei prodotti realizzati:

\begin{itemize}
    \item \textbf{\glo{RTB}{rtb}} (Requirements and Technology Baseline): Questa fase si concentra sull'analisi del problema, la definizione delle norme di lavoro, la stesura dei requisiti e lo studio delle tecnologie necessarie.
Include anche la realizzazione di un \textbf{Proof of Concept} (\textbf{PoC}) per validare la fattibilità dell'idea del prodotto e le scelte tecnologiche.
\item \textbf{\glo{PB}{pb}}/\textbf{\glo{MVP}{mvp}} (Product Baseline / Minimum Viable Product): In questa fase il focus si sposta sulla progettazione architetturale di dettaglio e sullo sviluppo effettivo del prodotto.
L'obiettivo è arrivare al rilascio di un \textit{MVP} funzionante validato, pronto per il collaudo formale con il \textbf{\glo{Committente}{committente}} e per l'approvazione da parte del \textbf{\glo{Proponente}{proponente} Var Group}.
\end{itemize}

Di seguito vengono dettagliate le sotto-\textit{milestone} per ciascuna macro-fase, con le relative tempistiche di previsione e l'elenco delle attività da svolgere.
(La tabella riassuntiva è riportata al termine del capitolo).

\subsection{Requirements and Technology Baseline (\textbf{RTB})}
La fase di lavoro dedicata al raggiungimento della \textit{milestone} \textbf{RTB} si estende dall'inizio di aprile fino a metà agosto ed è focalizzata sul consolidamento delle basi teoriche, metodologiche e tecnologiche del progetto. Inizialmente era previsto che si concludesse entro metà giugno. Seguono le sue sotto-\textit{milestone}, inizialmente definite:

\subsubsection{Fondamenta e normativa}
Previsione: Inizio Maggio \\
In questa fase iniziale il gruppo si concentra sull'organizzazione interna e sulla standardizzazione dei processi.
\begin{itemize}
    \item Definizione del \textit{Way of Working} del team.
    \item Redazione del documento \textbf{Norme di Progetto} (NdP).
\end{itemize}


\subsubsection{Analisi e studio}
Previsione: Metà Maggio \\
Fase dedicata alla comprensione dei bisogni del \textbf{Proponente} e alla definizione degli standard qualitativi.
\begin{itemize}
    \item \textbf{Analisi dei Requisiti (AdR)}: individuazione e definizione di tutti gli \textit{Use Case} e dei requisiti di sistema.
    \item \textbf{\glo{Piano di Qualifica (PdQ)}{piano-di-qualifica}}: definizione delle metriche di qualità.
    \item Completamento dello studio delle tecnologie necessarie per lo sviluppo dei tre agenti.
\end{itemize}

%%%%%%%%%%%%%%%%%%%%%%%

%%%%%%%%%%%%%%%%%%%%%%%%

\subsubsection{\textit{PoC}}
Previsione: Inizio Giugno \\
Fase pratica di realizzazione di alcuni requisiti precedentemente stabiliti e dell'implementazione delle scelte tecnologiche principali.
\begin{itemize}
    \item Pianificazione e realizzazione del \textit{Proof of Concept} (\textit{PoC}).
\end{itemize}

\subsubsection{Chiusura \textbf{RTB}}
Previsione: Metà Giugno \\
Consolidamento del lavoro svolto per la presentazione della prima \textit{milestone}.
\begin{itemize}
    \item Completamento di tutta la documentazione e verifica generale.
\item Preparazione ai colloqui di revisione \textbf{RTB}.
\end{itemize}


\subsubsection{Consuntivo della \textit{milestone} RTB}

Rispetto alle previsioni, la fase di lavoro dedicata alla preparazione di questa \textit{milestone} si è protratta per circa due mesi oltre quanto previsto. Ciò è dovuto a più ragioni:

\begin{itemize}
    \item definizione non sufficientemente adeguata e corretta di alcune sotto-\textit{milestone};

    \item implementazione e rispetto insufficiente delle sotto-\textit{milestone};

    \item tempi prolungati per l'adozione delle metodologie di lavoro definite nelle \textbf{Norme di Progetto};

    \item parte del lavoro svolto è risultata successivamente di qualità insufficiente e ha dovuto essere ripetuta.
\end{itemize}

Nella prossima fase di lavoro, il gruppo si impegnerà a prevenire il ripresentarsi degli stessi rischi. Sarà necessario aumentare l'efficienza del gruppo nel periodo rimanente, così da riuscire a concludere il progetto entro la scadenza finale del 25 settembre.


\subsection{Product Baseline (PB)}
La fase di sviluppo dedicata al raggiungimento della \textit{milestone} \textbf{PB} prende avvio a metà agosto e si conclude nell'ultima parte di settembre. L'obiettivo è tradurre il lavoro di analisi nello sviluppo del prodotto finale (\textit{MVP}). La fase è suddivisa nelle seguenti sotto-\textit{milestone}:

\subsubsection{Progettazione architetturale}
Previsione: \\
Passaggio dall'analisi alla progettazione.
\begin{itemize}
    \item Progettazione della struttura software sulla base delle tecnologie scelte, in modo da preparare il prodotto alla fase di codifica.

    \item Inizio della redazione del documento \textbf{Specifica Tecnica}.
\end{itemize}

%%%%%%%%%%%%%%%%%%%%%%%

\subsubsection{Sviluppo}
Previsione: \\
Fase di programmazione intensiva, focalizzata sul nucleo del sistema.
\begin{itemize}
    \item Codifica delle funzionalità classificate come obbligatorie.
    \item Codifica delle funzionalità associate ai requisiti non obbligatori, se le risorse disponibili lo permettono.
\end{itemize}

\subsubsection{Test e raffinamento}
Previsione: \\
Verifica e validazione del codice prodotto prima della presentazione al \textbf{Committente}.
\begin{itemize}
    \item Esecuzione dei \textit{Test di Sistema} e correzione dei \textit{bug} individuati.
    \item Revisione e aggiornamento della documentazione (v2.0.0).
\end{itemize}

\subsubsection{Validazione \textbf{MVP} e consegna}
Previsione: \\
Rilascio della prima versione funzionante del prodotto.
\begin{itemize}
    \item Redazione del \textbf{Manuale Utente} per facilitare l'utilizzo del sistema.
    \item Dimostrazione live dell'\textbf{MVP} prima al \textbf{Proponente} e poi al \textbf{Committente}, con relativa approvazione formale.
    \item Consegna di tutto il materiale, comprendente codice e documentazione.
\end{itemize}


\vspace{0.5cm}

\begin{table}[H]
    \centerline{
        \centering
        \renewcommand{\arraystretch}{1.3}
        % MODIFICA QUI: cambiato m{4em} in m{2.5cm} per dare più respiro al testo
        \begin{tabularx}{1.1\textwidth}{|m{5em}|p{11em}|p{7.5em}|X|}
            \hline
            \textbf{Macro \newline \textit{Milestone}} & \textbf{Sotto \textit{Milestone}} & \textbf{Previsione \newline completamento} & \textbf{Attività da svolgere} \\
            \hline
  
            %in multirow la misura tra le [] è l'offset verticale
            \multirow{4}{*}[-2cm]{\centering \textbf{RTB}} & Fondamenta e normativa & Inizio Maggio & - Definito \textit{Way of Working} \newline
            - NdP: redatto il documento \\
            \cline{2-4}
            & Analisi e ricerca & Metà Maggio & - \textbf{AdR}: Completati tutti gli \textit{Use Case} e requisiti \newline
            - \textbf{PdQ}: Definite le metriche \newline
            - Completato studio delle tecnologie \\ 
            \cline{2-4}
            &\textbf{PoC} & Inizio Giugno & - Pianificato e realizzato \textbf{PoC} \newline \\ 
       
            \cline{2-4}
            & Chiusura \textbf{RTB} & Metà Giugno & - Revisione e completamento di tutti i documenti (v1.0.0) \newline
            - Preparazione per i colloqui \textbf{RTB} \\
            \hline \hline
            \multirow{5}{*}[-3cm]{\centering \textbf{PB}} & Progettazione architetturale & 19 agosto & - Progettazione della struttura software basata sulle tecnologie scelte. \newline
            - Avvio della redazione della \textbf{Specifica Tecnica}. \\ 

            
            \cline{2-4}
            & Sviluppo & 31 agosto & - Codifica delle funzionalità obbligatorie. \newline
            - Codifica delle funzionalità non obbligatorie, se possibile. \\ 
            \cline{2-4}
            & Test e raffinamento & 11 settembre* & - Esecuzione dei \textit{Test di Sistema} e correzione dei \textit{bug}. \newline
            - Aggiornamento della documentazione alla versione 2.0.0. \\ 
            \cline{2-4}
            & Validazione \textbf{MVP} e consegna & 21 settembre & - Redazione del \textbf{Manuale Utente}. \newline
            - Dimostrazione live dell'\textbf{MVP} al \textbf{Proponente} e al \textbf{Committente}. \newline 
            - Consegna del codice e della documentazione. \\ 
            \cline{2-4}
  
            \hline
        \end{tabularx}
    }
\end{table}
* Aggiornato il 27/08.

\newpage
\section{\textit{Preventivo} e \textit{consuntivo}}

\subsection{Requirements and Technology Baseline (\textbf{RTB})}

\subsubsection{\textit{Sprint} 1: 1/04/2026 - 14/04/2026}

\paragraph{Obiettivi e attività}

Le attività previste per il primo \textit{sprint} sono le seguenti:

\begin{itemize}
\item
  Analizzare gli agenti proposti dall\textquotesingle azienda Var Group e
  decidere quali sviluppare per il progetto;
\item
  Strutturare ed iniziare la stesura dell\textquotesingle \textbf{Analisi dei Requisiti};
\item
  Strutturare i capitoli dei documenti: \textbf{Piano di Progetto},
  \textbf{Glossario} e \textbf{Norme di Progetto};
\item
  Introdurre le \textit{\glo{GitHub Actions}{github-actions}} per automatizzare
  l\textquotesingle aggiunta dei verbali sul sito web;
\item
  Comunicare all\textquotesingle azienda Var Group
  l\textquotesingle aggiudicazione del capitolato;
\item
  Studiare concetti e tecnologie necessarie per lo sviluppo del
  progetto.
\end{itemize}

A seguito di comunicazioni intercorse con l\textquotesingle azienda
\textbf{Proponente} Var Group, è emersa l\textquotesingle impossibilità da parte
dei referenti esterni di fornire supporto diretto fino alla metà di
aprile, a causa di scadenze interne.

Di conseguenza il gruppo si concentrerà maggiormente su attività che non
necessitano di un riscontro immediato da parte
dell\textquotesingle azienda.
\paragraph{\textit{Preventivo} orario ed economico}

I ruoli assegnati sono: \textbf{Responsabile}, \textbf{\glo{Amministratore}{amministratore}}, \textbf{\glo{Analista}{analista}} e \textbf{\glo{Verificatore}{verificatore}}.
\begin{table}[H]
    \centering
    \renewcommand{\arraystretch}{1.3}
    \begin{tabularx}{0.9\textwidth}{|X|c|c|c|c|c|c|c|}
        \hline
        \textbf{\textit{Sprint} 1} & \multicolumn{6}{|c|}{\textbf{Ruolo}} & \\ \hline
        \textbf{Membro} &\textbf{RE} & \textbf{AM} & \textbf{AN} & \textbf{PG} & \textbf{PR} & \textbf{VE} & \textbf{Ore totali} \\ \hline
        Andrea Fistarol & 2 & & & & & & \textbf{2} \\ \hline
        Dario Pirrone & & & 4 & & & & \textbf{4} \\ \hline
  
        Edoardo Granziero & & & & & & 2 & \textbf{2} \\ \hline
        Marco Barbiero & & 3 & & & & & \textbf{3} \\ \hline
        Samuele Bertin & & & 2 & & & & \textbf{2} \\ \hline
        Sara Ristovic & & & 4 & & & & \textbf{4} \\ \hline
        \textbf{Ore totali per ruolo} & \textbf{2} & \textbf{3} & \textbf{10} & \textbf{0} &
  
        \textbf{0} & \textbf{2} & \textbf{17} \\ \hline
    \end{tabularx}
    \caption{\textit{Preventivo} orario - \textit{sprint} 1}
\end{table}


Le risorse economiche previste per lo \textit{sprint} sono:

\begin{table}[H]
    \centering
    \renewcommand{\arraystretch}{1.3}
    \begin{tabularx}{0.9\textwidth}{|X|c|c|c|c|c|c|c|}
        \hline
        \textbf{\textit{Sprint} 1} & \textbf{RE} & \textbf{AM} & \textbf{AN} & \textbf{PG} & \textbf{PR} & \textbf{VE} & \textbf{Totale} \\ \hline
        \textbf{Costo orario} & \EUR{30} & \EUR{20} & \EUR{25} & \EUR{25} & \EUR{15} & \EUR{15} & / 
\\ \hline
        \textbf{Costo totale} & \EUR{60} & \EUR{60} & \EUR{250} & \EUR{0} & \EUR{0} & \EUR{30} & \textbf{\EUR{400}} \\ \hline
    \end{tabularx}
    \caption{\textit{Preventivo} economico - \textit{sprint} 1}
\end{table}

\paragraph{\textit{Consuntivo}}

È risultato un consumo orario minore di quello preventivato inizialmente
per \textbf{Amministratore} e \textbf{Verificatore}.
\begin{table}[H]
    \centering
    \renewcommand{\arraystretch}{1.3}
    \begin{tabularx}{0.9\textwidth}{|X|c|c|c|c|c|c|c|}
        \hline
        \textbf{\textit{Sprint} 1} & \multicolumn{6}{|c|}{\textbf{Ruolo}} & \\ \hline
        \textbf{Membro} & \textbf{RE} & \textbf{AM} & \textbf{AN} & \textbf{PG} & \textbf{PR} & \textbf{VE} & \textbf{Ore totali} \\ \hline
        Andrea Fistarol & 1 (-1) & & & & & & \textbf{1} \\ \hline
        Dario Pirrone & & & 4 & & & & \textbf{4} \\ 
\hline
        Edoardo Granziero & & & & & & 1 (-1) & \textbf{1} \\ \hline
        Marco Barbiero & & 3 & & & & & \textbf{3} \\ \hline
        Samuele Bertin & & & 2 & & & & \textbf{2} \\ \hline
        Sara Ristovic & & & 4 & & & & \textbf{4} \\ \hline
        \textbf{Ore totali per ruolo} & \textbf{1} & \textbf{3} & \textbf{10} & 
\textbf{0} & \textbf{0} & \textbf{1} & \textbf{15 (-2)} \\ \hline
    \end{tabularx}
    \caption{\textit{Consuntivo} orario - \textit{sprint} 1}
\end{table}

Le risorse economiche utilizzate durante lo \textit{sprint} sono:
\begin{table}[H]
    \centerline{
        \centering
        \renewcommand{\arraystretch}{1.3}
        \begin{tabularx}{1.09\textwidth}{|X|c|c|c|c|c|c|c|}
            \hline
            \textbf{\textit{Sprint} 1} & \textbf{RE} & \textbf{AM} & \textbf{AN} & \textbf{PG} & \textbf{PR} & \textbf{VE} & \textbf{Totale} \\ \hline
   
            \textbf{Costo orario} & \EUR{30} & \EUR{20} & \EUR{25} & \EUR{25} & \EUR{15} & \EUR{15} & / \\ \hline
            \textbf{Costo totale} & \EUR{30} & \EUR{60} & \EUR{250} & \EUR{0} & \EUR{0} & \EUR{15} & \textbf{\EUR{355}} \\ \hline
            \textbf{Saldo a fine \textit{sprint}} & \textbf{\EUR{1.590}} & \textbf{\EUR{900}} & \textbf{\EUR{1.400}} 
            & \textbf{\EUR{2.700}} & \textbf{\EUR{2.070}} & \textbf{\EUR{2.235}} & \textbf{\EUR{10.895}}\\ \hline
   
        \end{tabularx}
        }
    \caption{\textit{Consuntivo} economico - \textit{sprint} 1}
\end{table}

\paragraph{Analisi dei rischi occorsi}

Durante lo svolgimento del primo \textit{sprint}, si è manifestata una
criticità relativa alla definizione delle attività pianificate
(\textbf{RO1}): è stata riscontrata una definizione eccessivamente
generica di alcune attività (es.
"Inizio stesura \textbf{Analisi dei Requisiti}"). La mancanza di una scomposizione dettagliata ha causato
difficoltà nel monitoraggio delle attività da svolgere e nella gestione
del tempo da parte del gruppo.

Per mitigare questa criticità si è deciso di definire meglio il
\textit{Way of Working} riguardo a questo aspetto specifico.
\paragraph{\textit{\glo{Retrospettiva}{retrospettiva}}}

Al termine dello \textit{sprint} sono state portate a termine tutte le
attività previste, nonostante ciò il gruppo ha riscontrato
un\textquotesingle efficienza operativa inferiore al potenziale, la
causa principale è stata l\textquotesingle eccessiva genericità nella
definizione dei task iniziali (come esposto nell\textquotesingle analisi
dei rischi occorsi).

Per i prossimi \textit{sprint} il team ha deciso di integrare nel \textit{Way of Working} le seguenti migliorie:

\begin{itemize}
\item
  \textbf{Granularità dei \textit{Task}}: scomposizione delle
  macro-attività in \textit{\glo{Issue}{issue}} atomiche e ben definite in fase di
  \textit{\glo{Planning}{planning}}.
\item
  \textbf{Monitoraggio \textit{\glo{Data-driven}{data-driven}}}: Introduzione dei
  \textit{\glo{Burndown Chart}{burndown-chart}} per monitorare visivamente
  l\textquotesingle avanzamento dei lavori e la velocità di smaltimento
  delle attività residue.
\end{itemize}

\subsubsection{\textit{Sprint} 2: 15/04/2026 - 28/04/2026}

\paragraph{Obiettivi e attività}

Le attività previste per il secondo \textit{sprint} sono le seguenti:

\begin{itemize}
\item
  Creare un \textit{\glo{Google Sheets}{google-sheets}} per gestire la rotazione dei ruoli, per
  la contabilizzazione delle ore produttive di lavoro e per monitorare
  il budget utilizzato/rimanente;
\item
  Identificare le macro attività principali nello svolgimento del
  progetto;
\item
  \textbf{Analisi dei Requisiti}: scrivere gli \textit{Use Case} dei 3
  agenti da sviluppare;
\item
  \textbf{Norme di Progetto}: scrivere i processi primari e processi
  organizzativi;
\item
  \textbf{Piano di Progetto}: scrivere l\textquotesingle analisi dei
  rischi, \textit{Consuntivo} \textit{sprint} 1 e \textit{Preventivo} \textit{sprint} 2.
\end{itemize}

\paragraph{\textit{Preventivo} orario ed economico}

I ruoli assegnati sono: \textbf{Responsabile}, \textbf{Amministratore}, \textbf{Analista} e \textbf{Verificatore}.
\begin{table}[H]
    \centering
    \renewcommand{\arraystretch}{1.3}
    \begin{tabularx}{0.9\textwidth}{|X|c|c|c|c|c|c|c|}
        \hline
        \textbf{\textit{Sprint} 2} & \multicolumn{6}{|c|}{\textbf{Ruolo}} & \\ \hline
        \textbf{Membro} & \textbf{RE} & \textbf{AM} & \textbf{AN} &
        \textbf{PG} & \textbf{PR} & \textbf{VE} & \textbf{Totale} \\ \hline
        Andrea Fistarol & & & 2 & & & &\textbf{2} \\ \hline
        Dario Pirrone & 4 & & & & & &\textbf{4} \\ \hline
        Edoardo Granziero & & & 3 & & & &\textbf{3} \\ \hline
        Marco Barbiero & & & & & & 4 &\textbf{4} \\ \hline
        Samuele Bertin & & & 2 & & & &\textbf{2} \\ \hline
        Sara Ristovic & & 3 & & & & &\textbf{3} \\ \hline
        \textbf{Ore totali per ruolo} & \textbf{4} & \textbf{3} &
        \textbf{7} & \textbf{0} & \textbf{0} & \textbf{4} & \textbf{18}\\ \hline
 
    \end{tabularx}
    \caption{\textit{Preventivo} orario - \textit{sprint} 2}
\end{table}

Le risorse economiche previste per lo \textit{sprint} sono:

\begin{table}[H]
    \centering
    \renewcommand{\arraystretch}{1.3}
    \begin{tabularx}{0.9\textwidth}{|X|c|c|c|c|c|c|c|}
        \hline
        \textbf{\textit{Sprint} 2} & \textbf{RE} & \textbf{AM} & \textbf{AN} &
        \textbf{PG} & \textbf{PR} & \textbf{VE} & \textbf{Totale} \\ \hline
        \textbf{Costo orario} & \EUR{30} & \EUR{20} & \EUR{25} & \EUR{25} & \EUR{15}
        & \EUR{15} & 
/ \\ \hline
        \textbf{Costo totale} & \EUR{120} & \EUR{60} & \EUR{175} & \EUR{0} & \EUR{0} & 
        \EUR{60} & \textbf{\EUR{415}} \\ \hline
    \end{tabularx}
    \caption{\textit{Preventivo} economico - \textit{sprint} 2}
\end{table}

\paragraph{\textit{Consuntivo}}
È risultato un consumo orario maggiore di quello preventivato inizialmente per l'\textbf{Amministratore}.
\begin{table}[H]
    \centering
    \renewcommand{\arraystretch}{1.3}
    \begin{tabularx}{0.9\textwidth}{|X|c|c|c|c|c|c|c|}
        \hline
        \textbf{\textit{Sprint} 2} & \multicolumn{6}{|c|}{\textbf{Ruolo}} & \\ \hline
        \textbf{Membro} & \textbf{RE} & \textbf{AM} & \textbf{AN} &
        \textbf{PG} & \textbf{PR} & \textbf{VE} & \textbf{Totale} \\ \hline
        Andrea Fistarol & & & 2 & & & &\textbf{2} \\ \hline
        Dario Pirrone & 4 & & & & & &\textbf{4} \\ \hline
        Edoardo Granziero & & & 3 & & & &\textbf{3} \\ \hline
        Marco Barbiero & & & & & & 4 &\textbf{4} \\ \hline
        Samuele Bertin & & & 2 & & & &\textbf{2} \\ \hline
        Sara Ristovic & & 4 (+1) & & & & &\textbf{4} \\ \hline
        \textbf{Ore totali per ruolo} & \textbf{4} & \textbf{4} &
        \textbf{7} & \textbf{0} & \textbf{0} & \textbf{4} & \textbf{19 (+1)}\\ \hline
 
    \end{tabularx}
    \caption{\textit{Consuntivo} orario - \textit{sprint} 2}
\end{table}

Le risorse economiche utilizzate durante lo \textit{sprint} sono:

\begin{table}[H]
    \centerline{
        \centering
        \renewcommand{\arraystretch}{1.3}
        \begin{tabularx}{1.09\textwidth}{|X|c|c|c|c|c|c|c|}
            \hline
            \textbf{\textit{Sprint} 2} & \textbf{RE} & \textbf{AM} & \textbf{AN} & \textbf{PG} & \textbf{PR} & \textbf{VE} & \textbf{Totale} \\ \hline
            \textbf{Costo 
orario} & \EUR{30} & \EUR{20} & \EUR{25} & \EUR{25} & \EUR{15} & \EUR{15} & / \\ \hline
            \textbf{Costo totale} & \EUR{120} & \EUR{80} & \EUR{175} & \EUR{0} & \EUR{0} & \EUR{60} & \textbf{\EUR{435}} \\ \hline
            \textbf{Saldo a fine \textit{sprint}} & \textbf{\EUR{1.470}} & \textbf{\EUR{820}} & \textbf{\EUR{1.225}} 
            & \textbf{\EUR{2.700}} & \textbf{\EUR{2.070}} & \textbf{\EUR{2.175}} & \textbf{\EUR{10.460}}\\ \hline
        \end{tabularx}
    }
 
    \caption{\textit{Consuntivo} economico - \textit{sprint} 2}
\end{table}

\begin{figure}[H]
    \centering
    \includegraphics[width=0.9\textwidth]{grafici/burndown_sprint2.png}
    \caption{\textit{Grafico Burndown} - \textit{sprint} 2}
    \label{fig:burndown-sprint2}
\end{figure}

\paragraph{Analisi dei rischi occorsi}
Durante lo svolgimento del secondo \textit{sprint}, l'applicazione delle misure di mitigazione per il rischio \textbf{RO1 (Genericità dei task)} ha prodotto esiti positivi, tuttavia l'aumento della complessità gestionale ha fatto emergere una nuova criticità:

\begin{itemize}
    \item \textbf{Difficoltà nel coordinamento dei ruoli (RO2):} con l'introduzione di una scomposizione più fine delle attività in \textit{Issue} atomiche, è emersa una difficoltà iniziale nel coordinare il lavoro sincrono tra i diversi ruoli.
    \item \textbf{Mitigazione:} per risolvere tempestivamente l'intoppo, il gruppo ha deciso di avere una comunicazione più frequente tramite \textit{\glo{WhatsApp}{whatsapp}} e \textit{\glo{Discord}{discord}}. Questa misura ha permesso di ridistribuire il carico di lavoro in tempo reale, garantendo agli \textbf{Analisti} le linee guida necessarie per procedere senza ulteriori blocchi.
    \item \textbf{\textit{Issue} non chiuse:} non siamo riusciti a chiudere tutte le \textit{Issue} aperte a cause delle difficoltà che abbiamo incontrato nella gestione dei ruoli e delle proprie ore(come si vede nel grafico \textit{Burndown Chart}).
\end{itemize}

\paragraph{\textit{Retrospettiva}}
Il bilancio dello \textit{sprint} 2 mostra un netto miglioramento qualitativo nel metodo di lavoro e una gestione più consapevole delle risorse economiche e temporali rispetto al periodo precedente.

\paragraph{Cosa ha funzionato:}
\begin{itemize}
    \item \textbf{Efficacia della granularità:}La scomposizione delle macro-attività in \textit{Issue} atomiche ha permesso un monitoraggio molto più preciso;
    lo stato di avanzamento dell'\textbf{Analisi dei Requisiti} è risultato chiaro a tutti i membri del gruppo in ogni momento.
    \item \textbf{Monitoraggio \textit{Data-driven}:} l'adozione dei \textit{Burndown Chart} è stata fondamentale per identificare precocemente un rallentamento a metà \textit{sprint}, consentendo al \textbf{Responsabile} di intervenire proattivamente prima che la scadenza diventasse critica.
\end{itemize}

\paragraph{Migliorie per lo \textit{sprint} 3:}
\begin{itemize}
    \item \textbf{Gestione delle dipendenze:} nonostante l'atomicità dei task, è emersa la necessità di definire meglio l'ordine di esecuzione (propedeuticità) quando un'attività dipende da un'altra (es. la fase di verifica subordinata al completamento della stesura).
    \item \textbf{Standardizzazione dei documenti:} per ottimizzare i tempi di formattazione {\glo{\LaTeX{}}{latex}}, che sono risultati superiori al previsto, il team ha deciso di sviluppare dei \textit{template} pre-impostati per ogni tipologia di documento, velocizzando così la futura stesura tecnica.
\end{itemize}

\subsubsection{\textit{Sprint} 3: 29/04/2026 - 12/05/2026}

\paragraph{Obiettivi e attività}

Le attività previste per il terzo \textit{sprint} sono le seguenti:

\begin{itemize}
\item
    Iniziare il \textbf{Glossario};
\item
  Convertire i documenti in sviluppo (\textbf{Analisi dei Requisiti}, \textbf{Norme di Progetto} e \textbf{Piano di Progetto}) in \LaTeX{};
\item
  \textbf{Analisi dei Requisiti}: continuare gli \textit{Use Case} dei 3
  agenti da sviluppare;
\item
  \textbf{Norme di Progetto}: continuare i processi primari e processi
  organizzativi;
\item
  \textbf{Piano di Progetto}: stesura Pianificazione, scrivere \textit{Consuntivo} \textit{sprint} 2 e \textit{Preventivo} \textit{sprint} 3.
\end{itemize}


\paragraph{\textit{Preventivo} orario ed economico}

I ruoli assegnati sono: \textbf{Responsabile}, \textbf{Amministratore}, \textbf{Analista} e \textbf{Verificatore}.
\begin{table}[H]
    \centering
    \renewcommand{\arraystretch}{1.3}
    \begin{tabularx}{0.9\textwidth}{|X|c|c|c|c|c|c|c|}
        \hline
        \textbf{\textit{Sprint} 3} & \multicolumn{6}{|c|}{\textbf{Ruolo}} & \\ \hline
        \textbf{Membro} & \textbf{RE} & \textbf{AM} & \textbf{AN} &
        \textbf{PG} & \textbf{PR} & \textbf{VE} & \textbf{Totale} \\ \hline
        Andrea Fistarol & & 5 & & & & &\textbf{5} \\ \hline
        Dario Pirrone & & & & & & 3 &\textbf{3} \\ \hline
        Edoardo Granziero & & 4 & & & & &\textbf{4} \\ \hline
        Marco Barbiero & & & 3 & & & &\textbf{3} \\ \hline
        Samuele Bertin & 5 & & & & & &\textbf{5} \\ \hline
        Sara Ristovic & & & 8 & & & &\textbf{8} \\ \hline
        \textbf{Ore totali per ruolo} & \textbf{5} & \textbf{9} &
        \textbf{11} & \textbf{0} & \textbf{0} & \textbf{3} & \textbf{28}\\ \hline
 
    \end{tabularx}
    \caption{\textit{Preventivo} orario - \textit{sprint} 3}
\end{table}

Le risorse economiche previste per lo \textit{sprint} sono:

\begin{table}[H]
    \centering
    \renewcommand{\arraystretch}{1.3}
    \begin{tabularx}{0.9\textwidth}{|X|c|c|c|c|c|c|c|}
        \hline
        \textbf{\textit{Sprint} 3} & \textbf{RE} & \textbf{AM} & \textbf{AN} &
        \textbf{PG} & \textbf{PR} & \textbf{VE} & \textbf{Totale} \\ \hline
        \textbf{Costo orario} & \EUR{30} & \EUR{20} & \EUR{25} & \EUR{25} & \EUR{15}
        & \EUR{15} & 
/ \\ \hline
        \textbf{Costo totale} & \EUR{150} & \EUR{180} & \EUR{275} & \EUR{0} & \EUR{0} & 
        \EUR{45} & \textbf{\EUR{650}} \\ \hline
    \end{tabularx}
    \caption{\textit{Preventivo} economico - \textit{sprint} 3}
\end{table}

\paragraph{\textit{Consuntivo}}

È risultato un consumo orario in linea con quanto preventivato inizialmente per tutti i ruoli.
\begin{table}[H]
    \centering
    \renewcommand{\arraystretch}{1.3}
    \begin{tabularx}{0.9\textwidth}{|X|c|c|c|c|c|c|c|}
        \hline
        \textbf{\textit{Sprint} 3} & \multicolumn{6}{|c|}{\textbf{Ruolo}} & \\ \hline
        \textbf{Membro} & \textbf{RE} & \textbf{AM} & \textbf{AN} &
        \textbf{PG} & \textbf{PR} & \textbf{VE} & \textbf{Totale} \\ \hline
        Andrea Fistarol & & 5 & & & & &\textbf{5} \\ \hline
        Dario Pirrone & & & & & & 3 &\textbf{3} \\ \hline
        Edoardo Granziero & & 4 & & & & &\textbf{4} \\ \hline
        Marco Barbiero & & & 3 & & & &\textbf{3} \\ \hline
        Samuele Bertin & 5 & & & & & &\textbf{5} \\ \hline
        Sara Ristovic & & & 8 & & & &\textbf{8} \\ \hline
        \textbf{Ore totali per ruolo} & \textbf{5} & \textbf{9} &
        \textbf{11} & \textbf{0} & \textbf{0} & \textbf{3} & \textbf{28}\\ \hline
 
    \end{tabularx}
    \caption{\textit{Consuntivo} orario - \textit{sprint} 3}
\end{table}

Le risorse economiche utilizzate durante lo \textit{sprint} sono:

\begin{table}[H]
    \centerline{
        \centering
        \renewcommand{\arraystretch}{1.3}
        \begin{tabularx}{1.09\textwidth}{|X|c|c|c|c|c|c|c|}
            \hline
            \textbf{\textit{Sprint} 3} & \textbf{RE} & \textbf{AM} & \textbf{AN} &
            \textbf{PG} & \textbf{PR} & \textbf{VE} & \textbf{Totale} \\ \hline
  
            \textbf{Costo orario} & \EUR{30} & \EUR{20} & \EUR{25} & \EUR{25} & \EUR{15}
            & \EUR{15} & / \\ \hline
            \textbf{Costo totale} & \EUR{150} & \EUR{180} & \EUR{275} & \EUR{0} & \EUR{0} & 
            \EUR{45} & \textbf{\EUR{650}} \\ \hline
            \textbf{Saldo a fine \textit{sprint}} & \textbf{\EUR{1.320}} & \textbf{\EUR{640}} & 
\textbf{\EUR{950}} 
            & \textbf{\EUR{2.700}} & \textbf{\EUR{2.070}} & \textbf{\EUR{2.130}} & \textbf{\EUR{9.810}}\\ \hline
        \end{tabularx}
    }
    \caption{\textit{Consuntivo} economico - \textit{sprint} 3}
\end{table}

\begin{figure}[H]
    \centering
    \includegraphics[width=0.9\textwidth]{grafici/burndown_sprint3.png}
    \caption{\textit{Grafico Burndown} - \textit{sprint} 3}
    \label{fig:burndown-sprint3}
\end{figure}

\paragraph{Analisi dei rischi occorsi}

Durante lo svolgimento del terzo \textit{sprint} è riemersa una criticità legata alla pianificazione (\textbf{RO1 - Scarsa Pianificazione}).
La divisione delle \textit{Issue} non si è rivelata sufficientemente granulare: avere task troppo ampi non ci ha permesso di tracciare i progressi nel migliore dei modi.
Questo ha reso la lettura del \textit{Burndown Chart} poco indicativa dello stato reale dei lavori e, soprattutto, ha generato un collo di bottiglia, impedendoci di chiudere definitivamente diverse attività prima della fine dello \textit{sprint}.
\paragraph{\textit{Retrospettiva}}

Il terzo \textit{sprint} ha consolidato l'infrastruttura documentale, ma ha evidenziato evidenti limiti operativi nella gestione e chiusura dei task.

\textbf{Cosa ha funzionato:}
\begin{itemize}
    \item \textbf{Standardizzazione in }\LaTeX{}: l'adozione dei nuovi \textit{template} in \LaTeX{} ha reso il lavoro documentale molto più ordinato.
    L'omogeneità stilistica e la struttura preimpostata hanno facilitato notevolmente la stesura, ripagando l'investimento di tempo fatto nello \textit{sprint} precedente e rendendo i documenti molto più manutenibili.
\end{itemize}

\textbf{Cosa migliorare per lo \textit{sprint} 4:}
\begin{itemize}
    \item \textbf{Scomposizione atomica e chiusura delle \textit{Issue}:} a causa della scarsa granularità, molte attività previste non sono arrivate allo stato di "Done" e sono slittate allo \textit{sprint} successivo.
    Per lo \textit{sprint} 4, il team si impegnerà a spezzare i task in unità di lavoro molto più piccole e indipendenti, per garantirne la chiusura tempestiva, evitare lo \textit{\glo{spillover}{spillover}} tra \textit{sprint} e permettere un tracciamento continuo e reale dei progressi.
\end{itemize}

\subsubsection{\textit{Sprint} 4: 13/05/2026 - 26/05/2026}

\paragraph{Obiettivi e attività}

Le attività previste per il quarto \textit{sprint} sono le seguenti:

\begin{itemize}
    \item \textbf{Analisi dei Requisiti}: completamento degli \textit{Use Case} per gli agenti \textit{README}, \textit{Changelog} e \textit{OWASP}, stesura della sezione di Tracciamento e completamento della tabella dei Requisiti;
\item \textbf{Piano di Qualifica}: stesura degli obiettivi di qualità di processo;
\item \textbf{Norme di Progetto}: stesura del capitolo relativo allo sviluppo;
\item \emph{Gestione documentale e \textit{Repository}:} aggiornamento del \textbf{Glossario} con introduzione di nuovi termini e creazione dei collegamenti tra documenti e \textbf{Glossario}, sistemazione della struttura delle cartelle \textit{\glo{GitHub}{github}} per i file \texttt{.tex} e sistemazione dei \textit{template} dei documenti;
\end{itemize}

\paragraph{\textit{Preventivo} orario ed economico}

I ruoli assegnati sono: \textbf{Responsabile}, \textbf{Amministratore}, \textbf{Analista} e \textbf{Verificatore}.
\begin{table}[H]
    \centering
    \renewcommand{\arraystretch}{1.3}
    \begin{tabularx}{0.9\textwidth}{|X|c|c|c|c|c|c|c|}
        \hline
        \textbf{\textit{Sprint} 4} & \multicolumn{6}{|c|}{\textbf{Ruolo}} & \\ \hline
        \textbf{Membro} & \textbf{RE} & \textbf{AM} & \textbf{AN} &
        \textbf{PG} & \textbf{PR} & \textbf{VE} & \textbf{Totale} \\ \hline
        Andrea Fistarol & & & 4 & & & & \textbf{4} \\ \hline
        Dario Pirrone & & 5 & & & & & \textbf{5} \\ \hline
        Edoardo Granziero & 4 & & & & & & \textbf{4} \\ \hline
        Marco Barbiero & & & 5 & & & & \textbf{5} \\ \hline
        Samuele Bertin & & & & & & 3 & \textbf{3} \\ \hline
        Sara Ristovic & & & 6 & & & & \textbf{6} \\ \hline
        \textbf{Ore totali per ruolo} & \textbf{4} & \textbf{5} &
        \textbf{15} & \textbf{0} & \textbf{0} & \textbf{3} & \textbf{27}\\ \hline
    \end{tabularx}
    \caption{\textit{Preventivo} orario - \textit{sprint} 4}
\end{table}

Le risorse economiche previste per lo \textit{sprint} sono:

\begin{table}[H]
    \centering
    \renewcommand{\arraystretch}{1.3}
    \begin{tabularx}{0.9\textwidth}{|X|c|c|c|c|c|c|c|}
        \hline
        \textbf{\textit{Sprint} 4} & \textbf{RE} & \textbf{AM} & \textbf{AN} &
        \textbf{PG} & \textbf{PR} & \textbf{VE} & \textbf{Totale} \\ \hline
        \textbf{Costo orario} & \EUR{30} & \EUR{20} & \EUR{25} & \EUR{25} & \EUR{15}
     
        & \EUR{15} & / \\ \hline
        \textbf{Costo totale} & \EUR{120} & \EUR{100} & \EUR{375} & \EUR{0} & \EUR{0} & \EUR{45} & \textbf{\EUR{640}} \\ \hline
    \end{tabularx}
    \caption{\textit{Preventivo} economico - \textit{sprint} 4}
\end{table}

\paragraph{\textit{Consuntivo}}

È risultato un consumo orario leggermente superiore a quanto preventivato inizialmente.

\begin{table}[H]
    \centering
    \renewcommand{\arraystretch}{1.3}
    \begin{tabularx}{0.95\textwidth}{|X|c|c|c|c|c|c|c|}
        \hline
        \textbf{\textit{Sprint} 4} & \multicolumn{6}{|c|}{\textbf{Ruolo}} & \\ \hline
        \textbf{Membro} & \textbf{RE} & \textbf{AM} & \textbf{AN} &
        \textbf{PG} & \textbf{PR} & \textbf{VE} & \textbf{Totale} \\ \hline
        Andrea Fistarol & & & 4 & & & &\textbf{4} \\ \hline
        Dario Pirrone & & 4 (-1) & & & & &\textbf{4} \\ \hline
        Edoardo Granziero & 4 & & & & & &\textbf{4} \\ \hline
        Marco Barbiero & & & 5 & & & &\textbf{5} \\ \hline
        Samuele Bertin & & & & & & 4 (+1) &\textbf{4} \\ \hline
        Sara Ristovic & & & 8 (+2) & & & &\textbf{8} \\ \hline
        \textbf{Ore totali per ruolo} & \textbf{4} & \textbf{4} &
        \textbf{17} & \textbf{0} & \textbf{0} & \textbf{4} & \textbf{29 (+2)}\\ \hline
 
    \end{tabularx}
    \caption{\textit{Consuntivo} orario - \textit{sprint} 4}
\end{table}

Le risorse economiche utilizzate durante lo \textit{sprint} sono:

\begin{table}[H]
    \centerline{
        \centering
        \renewcommand{\arraystretch}{1.3}
        \begin{tabularx}{1.09\textwidth}{|X|c|c|c|c|c|c|c|}
            \hline
            \textbf{\textit{Sprint} 4} & \textbf{RE} & \textbf{AM} & \textbf{AN} &
            \textbf{PG} & \textbf{PR} & \textbf{VE} & \textbf{Totale} \\ \hline
  
            \textbf{Costo orario} & \EUR{30} & \EUR{20} & \EUR{25} & \EUR{25} & \EUR{15}
            & \EUR{15} & / \\ \hline
            \textbf{Costo totale} & \EUR{120} & \EUR{80} & \EUR{425} & \EUR{0} & \EUR{0} & 
            \EUR{60} & \textbf{\EUR{685}} \\ \hline
            \textbf{Saldo a fine \textit{sprint}} & \textbf{\EUR{1.200}} & \textbf{\EUR{560}} & \textbf{\EUR{525}} & \textbf{\EUR{2.700}} & \textbf{\EUR{2.070}} & \textbf{\EUR{2.070}} & \textbf{\EUR{9.125}}\\ \hline
        \end{tabularx}
    }
    \caption{\textit{Consuntivo} economico - \textit{sprint} 4}
\end{table}

\begin{figure}[H]
    \centering
    \includegraphics[width=0.9\textwidth]{grafici/burndown_sprint4.png}
    \caption{\textit{Grafico Burndown} - \textit{sprint} 4}
    \label{fig:burndown-sprint4}
\end{figure}

\paragraph{Analisi dei rischi occorsi}
Durante lo \textit{sprint 4} si è manifestato il rischio relativo alla comunicazione inefficiente con il \textbf{Proponente}, identificato come \textbf{RO3}. \\
A causa di tale difficoltà comunicativa, è stato necessario attendere più del previsto per fissare un incontro con il \textbf{Proponente}, finalizzato alla discussione dell'\textbf{Analisi dei Requisiti} e alla ricezione del primo riscontro significativo sul documento. Ciò ha causato un ritardo nello svolgimento del progetto e la necessità di rivalutare alcune scadenze interne.

\paragraph{\textit{Retrospettiva}}
Nonostante il rischio occorso, è stato comunque effettuato un buon avanzamento nella stesura di tutti i documenti previsti finora.

\paragraph{Cosa ha funzionato:}
\begin{itemize}
    \item \textbf{\textit{Issue} "non assegnate":} questo tipo di compito è stato introdotto per garantire maggiore flessibilità ai membri del gruppo. Qualora un membro avesse completato le proprie \textit{issue} assegnate prima della fine dello \textit{sprint} e avesse avuto ancora tempo a disposizione, poteva assegnarsi una delle \textit{issue} inizialmente indicate come non assegnate.  Tutte le \textit{issue} di questo tipo sono state completate, contribuendo ad aumentare la produttività del gruppo.
    \item \textbf{Comunicazione interna sullo stato di avanzamento a metà dello \textit{sprint}:} il gruppo ha ritenuto necessario comunicare più spesso lo stato di avanzamento del lavoro di ogni membro durante lo \textit{sprint}, poiché gli \textit{sprint} durano due settimane e si vuole promuovere un lavoro continuo. È stato concluso che questa comunicazione asincrona sullo stato di avanzamento ha aiutato i membri a lavorare meglio, dato che lo \textit{sprint} si è concluso con tutte le \textit{issue} chiuse, risolvendo una criticità emersa nello \textit{sprint} precedente.
\end{itemize}

\paragraph{Migliorie per lo \textit{sprint} 5:}
\begin{itemize}
    \item \textbf{Mitigazioni per il rischio riscontrato:} in futuro, durante la comunicazione tramite messaggi su \textit{Slack}, il gruppo taggherà sempre la persona esterna a cui si rivolge, poiché questa pratica si è dimostrata utile. Inoltre, le richieste importanti da parte del gruppo verranno formulate anche in forma scritta e non solamente oralmente durante gli incontri esterni.
    \item \textbf{Incontro interno a metà dello \textit{sprint}:} 
    si è deciso di svolgere questo incontro per monitorare meglio l’avanzamento dei compiti dello \textit{sprint}, pratica che si è rivelata molto utile nello \textit{sprint} precedente, e per concordare più aspetti importanti.
\end{itemize}


\subsubsection{\textit{Sprint} 5: 27/05/2026 - 09/06/2026}

\paragraph{Obiettivi e attività}

Le attività previste per il quinto \textit{sprint}, ordinate in base alla priorità, sono le seguenti:

\begin{itemize}
    \item Ricevere riscontro sull'\textbf{Analisi dei Requisiti}, effettuare le modifiche necessarie e introdurre i \textit{diagrammi UML} degli \textit{use cases}
    \item Decidere le tecnologie per l'implementazione del \textbf{PoC} e, se possibile, iniziare la sua realizzazione
    \item Avanzare nella stesura del \textbf{Piano di Qualifica} con la sezione relativa ai test
\end{itemize}

\paragraph{\textit{Preventivo} orario ed economico}

I ruoli assegnati sono: \textbf{Responsabile}, \textbf{Amministratore}, \textbf{Analista}, \textbf{Progettista} e \textbf{Verificatore}.
\begin{table}[H]
    \centering
    \renewcommand{\arraystretch}{1.3}
    \begin{tabularx}{0.9\textwidth}{|X|c|c|c|c|c|c|c|}
        \hline
        \textbf{\textit{Sprint} 5} & \multicolumn{6}{|c|}{\textbf{Ruolo}} & \\ \hline
        \textbf{Membro} & \textbf{RE} & \textbf{AM} & \textbf{AN} &
        \textbf{PG} & \textbf{PR} & \textbf{VE} & \textbf{Totale} \\ \hline
        Andrea Fistarol & & & 5 & & & & \textbf{5} \\ \hline
        Dario Pirrone & & & 7 & & & & \textbf{7} \\ \hline
        Edoardo Granziero & & & & & & 3 & \textbf{3} \\ \hline
        Marco Barbiero & & & & 8 & & & \textbf{8} \\ \hline
        Samuele Bertin & & 5 & & & & & \textbf{5} \\ \hline
        Sara Ristovic & 5 & & & & & & \textbf{5} \\ \hline
        \textbf{Ore totali per ruolo} & \textbf{5} & \textbf{5} &
        \textbf{12} & \textbf{8} & \textbf{0} & \textbf{3} & \textbf{33}\\ \hline
    \end{tabularx}
    \caption{\textit{Preventivo} orario - \textit{sprint} 5}
\end{table}

Le risorse economiche previste per lo \textit{sprint} sono:

\begin{table}[H]
    \centering
    \renewcommand{\arraystretch}{1.3}
    \begin{tabularx}{0.9\textwidth}{|X|c|c|c|c|c|c|c|}
        \hline
        \textbf{\textit{Sprint} 5} & \textbf{RE} & \textbf{AM} & \textbf{AN} &
        \textbf{PG} & \textbf{PR} & \textbf{VE} & \textbf{Totale} \\ \hline
        \textbf{Costo orario} & \EUR{30} & \EUR{20} & \EUR{25} & \EUR{25} & \EUR{15}
     
        & \EUR{15} & / \\ \hline
        \textbf{Costo totale} & \EUR{150} & \EUR{100} & \EUR{300} & \EUR{200} & \EUR{0} & \EUR{45} & \textbf{\EUR{795}} \\ \hline
    \end{tabularx}
    \caption{\textit{Preventivo} economico - \textit{sprint} 5}
\end{table}

\paragraph{\textit{Consuntivo}}

È risultato un consumo orario leggermente superiore a quanto  preventivato inizialmente.

\begin{table}[H]
    \centering
    \renewcommand{\arraystretch}{1.3}
    \begin{tabularx}{0.95\textwidth}{|X|c|c|c|c|c|c|c|}
        \hline
        \textbf{\textit{Sprint} 5} & \multicolumn{6}{|c|}{\textbf{Ruolo}} & \\ \hline
        \textbf{Membro} & \textbf{RE} & \textbf{AM} & \textbf{AN} &
        \textbf{PG} & \textbf{PR} & \textbf{VE} & \textbf{Totale} \\ \hline
        Andrea Fistarol & & & 5 & & & & \textbf{5} \\ \hline
        Dario Pirrone & & & 8 (+1) & & & & \textbf{8} \\ \hline
        Edoardo Granziero & & & & & & 4 (+1) & \textbf{4} \\ \hline
        Marco Barbiero & & & & 8 & & & \textbf{8} \\ \hline
        Samuele Bertin & & 4 (-1) & & & & & \textbf{4} \\ \hline
        Sara Ristovic & 5 & & & & & & \textbf{5} \\ \hline
        \textbf{Ore totali per ruolo} & \textbf{5} & \textbf{4} &
        \textbf{13} & \textbf{8} & \textbf{0} & \textbf{4} & \textbf{34 (+1)}\\ \hline
 
    \end{tabularx}
    \caption{\textit{Consuntivo} orario - \textit{sprint} 5}
\end{table}

Le risorse economiche utilizzate durante lo \textit{sprint} sono:

\begin{table}[H]
    \centerline{
        \centering
        \renewcommand{\arraystretch}{1.3}
        \begin{tabularx}{1.09\textwidth}{|X|c|c|c|c|c|c|c|}
            \hline
            \textbf{\textit{Sprint} 5} & \textbf{RE} & \textbf{AM} & \textbf{AN} &
            \textbf{PG} & \textbf{PR} & \textbf{VE} & \textbf{Totale} \\ \hline
  
            \textbf{Costo orario} & \EUR{30} & \EUR{20} & \EUR{25} & \EUR{25} & \EUR{15}
            & \EUR{15} & / \\ \hline
            \textbf{Costo totale} & \EUR{150} & \EUR{80} & \EUR{325} & \EUR{200} & \EUR{0} & 
            \EUR{60} & \textbf{\EUR{815}} \\ \hline
            \textbf{Saldo a fine \textit{sprint}} & \textbf{\EUR{1.050}} & \textbf{\EUR{480}} & \textbf{\EUR{200}} & \textbf{\EUR{2.500}} & \textbf{\EUR{2.070}} & \textbf{\EUR{2.010}} & \textbf{\EUR{8.310}}\\ \hline
        \end{tabularx}
    }
    \caption{\textit{Consuntivo} economico - \textit{sprint} 5}
\end{table}

\begin{figure}[H]
    \centering
    \includegraphics[width=0.9\textwidth]{grafici/burndown_sprint5.png}
    \caption{\textit{Grafico Burndown} - \textit{sprint} 5}
    \label{fig:burndown-sprint5}
\end{figure}

\paragraph{Analisi dei rischi occorsi}
\begin{itemize}
     \item Durante lo \textit{sprint} si è concretizzato il rischio di un disallineamento informativo con il \textbf{Proponente} riguardo l'architettura di partenza. L'assenza di un \textbf{orchestratore} agentico preesistente sviluppato da \textbf{Var Group}, contrariamente a quanto dichiarato precedentemente dall'azienda, richiederà uno sforzo aggiuntivo.\\
     L'imprevisto ha causato un incremento del carico di lavoro e un rallentamento minore sulla tabella di marcia originaria, portando alla necessità di un cambiamento all'\textbf{Analisi dei Requisiti}, questa dinamica è riconducibile al rischio \textbf{RR1}.
\end{itemize}

\paragraph{\textit{Retrospettiva}} In questo sprint il team è stato in grado di gestire i rischi occorsi senza grandi rallenamenti e riuscendo a chiudere tutte le \textit{issue} preventivate ad inizio sprint.
\paragraph{Cosa ha funzionato:}
\begin{itemize}
    \item \textbf{Avvio tempestivo e strutturato del \textbf{PoC}:} nonostante l'imprevisto tecnico riscontrato con l'azienda, l'assegnazione mirata delle ore al ruolo di \textit{Progettista} ha permesso di definire con successo lo scheletro architetturale del \textit{Proof of Concept} e di completare una parte significativa del suo sviluppo funzionale.
    \item \textbf{Raffinamento proattivo dei casi d'uso e requisiti:} gli \textit{Analisti} hanno individuato le aree critiche dell'\textbf{Analisi dei Requisiti} e hanno avviato una sessione di correzione dei \textit{casi d'uso} e irrobustimento del tracciamento, preparandosi al colloquio di revisione con il Prof. Cardin.
\end{itemize}

\paragraph{Migliorie per lo \textit{sprint} 6}
\begin{itemize}
    \item \textbf{Parallelizzazione tra sviluppo del PoC e correzione documentale:} sfruttando l'insieme delle tecnologie del \textbf{PoC} ormai consolidato, si prevede di separare nettamente i flussi di lavoro nello \textit{sprint} 6, che sarà caratterizzato da un calo della produttività e comunicazione intera a causa della concomitanza con la sessioe di esami.
\end{itemize}

\subsubsection{\textit{Sprint} 6: 10/06/2026 - 23/06/2026}

\paragraph{Obiettivi e attività}

Le attività previste per il sesto \textit{sprint}, ordinate in base alla priorità, sono le seguenti:

\begin{itemize}
    \item  Modifiche all'\textbf{Analisi dei Requisiti} in preparzione del ricevimento con il Prof. Cardin
    \item Procedere con la codifica del \textbf{PoC}
    \item Controllo approfondito di quanto definito nel \textbf{PdQ}
\end{itemize}

\paragraph{\textit{Preventivo} orario ed economico}

I ruoli assegnati sono: \textbf{Responsabile}, \textbf{Amministratore}, \textbf{Analista}, \textbf{Progettista} e \textbf{Verificatore}.
\begin{table}[H]
    \centering
    \renewcommand{\arraystretch}{1.3}
    \begin{tabularx}{0.9\textwidth}{|X|c|c|c|c|c|c|c|}
        \hline
        \textbf{\textit{Sprint} 6} & \multicolumn{6}{|c|}{\textbf{Ruolo}} & \\ \hline
        \textbf{Membro} & \textbf{RE} & \textbf{AM} & \textbf{AN} &
        \textbf{PG} & \textbf{PR} & \textbf{VE} & \textbf{Totale} \\ \hline
        Andrea Fistarol & & & & 1 & & &\textbf{1} \\ \hline
        Dario Pirrone & & & 2 & & & & \textbf{2} \\ \hline
        Edoardo Granziero & & 1 & & & & & \textbf{1} \\ \hline
        Marco Barbiero & & & & 1 & & & \textbf{1} \\ \hline
        Samuele Bertin & 2 & & & & & & \textbf{2} \\ \hline
        Sara Ristovic & & & & & & 1 & \textbf{1} \\ \hline
        \textbf{Ore totali per ruolo} & \textbf{2} & \textbf{1} &
        \textbf{2} & \textbf{2} & \textbf{0} & \textbf{1} & \textbf{8}\\ \hline
    \end{tabularx}
    \caption{\textit{Preventivo} orario - \textit{sprint} 6}
\end{table}

Le risorse economiche previste per lo \textit{sprint} sono:

\begin{table}[H]
    \centering
    \renewcommand{\arraystretch}{1.3}
    \begin{tabularx}{0.9\textwidth}{|X|c|c|c|c|c|c|c|}
        \hline
        \textbf{\textit{Sprint} 6} & \textbf{RE} & \textbf{AM} & \textbf{AN} &
        \textbf{PG} & \textbf{PR} & \textbf{VE} & \textbf{Totale} \\ \hline
        \textbf{Costo orario} & \EUR{30} & \EUR{20} & \EUR{25} & \EUR{25} & \EUR{15}
     
        & \EUR{15} & / \\ \hline
        \textbf{Costo totale} & \EUR{60} & \EUR{20} & \EUR{50} & \EUR{50} & \EUR{0} & 
            \EUR{15} & \textbf{\EUR{195}} \\ \hline
    \end{tabularx}
    \caption{\textit{Preventivo} economico - \textit{sprint} 6}
\end{table}

\paragraph{\textit{Consuntivo}}

È risultato un consumo orario in linea con quanto preventivato inizialmente.

\begin{table}[H]
    \centering
    \renewcommand{\arraystretch}{1.3}
    \begin{tabularx}{0.9\textwidth}{|X|c|c|c|c|c|c|c|}
        \hline
        \textbf{\textit{Sprint} 6} & \multicolumn{6}{|c|}{\textbf{Ruolo}} & \\ \hline
        \textbf{Membro} & \textbf{RE} & \textbf{AM} & \textbf{AN} &
        \textbf{PG} & \textbf{PR} & \textbf{VE} & \textbf{Totale} \\ \hline
        Andrea Fistarol & & & & 1 & & & \textbf{1} \\ \hline
        Dario Pirrone & & & 2 & & & & \textbf{2} \\ \hline
        Edoardo Granziero & & 1 & & & & & \textbf{1} \\ \hline
        Marco Barbiero & & & & 1 & & & \textbf{1} \\ \hline
        Samuele Bertin & 2 & & & & & & \textbf{2} \\ \hline
        Sara Ristovic & & & & & & 1 & \textbf{1} \\ \hline
        \textbf{Ore totali per ruolo} & \textbf{2} & \textbf{1} &
        \textbf{2} & \textbf{2} & \textbf{0} & \textbf{1} & \textbf{8}\\ \hline
 
    \end{tabularx}
    \caption{\textit{Consuntivo} orario - \textit{sprint} 6}
\end{table}

Le risorse economiche utilizzate durante lo \textit{sprint} sono:

\begin{table}[H]
    \centerline{
        \centering
        \renewcommand{\arraystretch}{1.3}
        \begin{tabularx}{1.09\textwidth}{|X|c|c|c|c|c|c|c|}
            \hline
            \textbf{\textit{Sprint} 6} & \textbf{RE} & \textbf{AM} & \textbf{AN} &
            \textbf{PG} & \textbf{PR} & \textbf{VE} & \textbf{Totale} \\ \hline
  
            \textbf{Costo orario} & \EUR{30} & \EUR{20} & \EUR{25} & \EUR{25} & \EUR{15}
            & \EUR{15} & / \\ \hline
            \textbf{Costo totale} & \EUR{60} & \EUR{20} & \EUR{50} & \EUR{50} & \EUR{0} & 
            \EUR{15} & \textbf{\EUR{195}} \\ \hline
            \textbf{Saldo a fine \textit{sprint}} & \textbf{\EUR{990}} & \textbf{\EUR{460}} & \textbf{\EUR{150}} & \textbf{\EUR{2.450}} & \textbf{\EUR{2.070}} & \textbf{\EUR{1.995}} & \textbf{\EUR{8.115}}\\ \hline
        \end{tabularx}
    }
    \caption{\textit{Consuntivo} economico - \textit{sprint} 6}
\end{table}

\begin{figure}[H]
  \centering
    \includegraphics[width=0.9\textwidth]{grafici/burndown_sprint6.png}
    \caption{\textit{Grafico Burndown} - \textit{sprint} 6}
    \label{fig:burndown-sprint6}
\end{figure}

\paragraph{Analisi dei rischi occorsi}
\begin{itemize}
     \item Durante lo \textit{sprint} 6 si è concretizzato il rischio legato alla concomitanza con la sessione d'esami estiva, che ha comportato una drastica e prevista riduzione della disponibilità oraria dei membri del team, riconducibile al rischio \textbf{RI1}. Questo vincolo accademico personale ha rallentato temporaneamente il ritmo delle attività di sviluppo e stesura documentale.\\
\end{itemize}

\paragraph{\textit{Retrospettiva}}Questo sprint è stato caratterizzato da una produttività ridotta a causa della concomitanza con la sessione esami, questa dinamica era stata ampiamente preventivata e il team è riuscito a fare i progressi deisiderati.
\paragraph{Cosa ha funzionato:}
\begin{itemize}
    \item \textbf{Gestione del carico di lavoro minimo:} nonostante le ore complessive dedicate al progetto siano state sensibilmente inferiori rispetto ai periodi precedenti, la pianificazione interna ha permesso di mantenere attivi i canali di comunicazione stabili e di non bloccare del tutto le attività propedeutiche.
    \item \textbf{Mantenimento dei requisiti di base:} il team è riuscito a preservare l'integrità e il consolidamento di quanto sviluppato finora sul \textbf{PoC} e sull'\textbf{Analisi dei Requisiti} durante gli \textit{sprint} passati, evitando regressioni o perdite di focus architetturale nonostante la ridotta operatività.
\end{itemize}

\paragraph{Migliorie per lo \textit{sprint} 7}
\begin{itemize}
    \item \textbf{Pianificazione di recupero e riallocazione del budget orario:} con la conclusione della parte principale della sessione d'esami, lo \textit{sprint} 7 dovrà prevedere una rimodulazione intensiva delle ore per recuperare il disavanzo dello \textit{sprint} precedente, parallelizzando lo sviluppo pratico del \textbf{PoC} e la chiusura dei documenti.
\end{itemize}

\subsubsection{\textit{Sprint} 7: 24/06/2026 - 07/07/2026}

\paragraph{Obiettivi e attività}

Le attività previste per il settimo \textit{sprint}, ordinate in base alla priorità, sono le seguenti:

\begin{itemize}
    \item  Verifica di tutta la documentazione in vista dell'\textbf{RTB}.
    \item Sviluppo \textbf{\glo{Frontend}{frontend}}, \textbf{\glo{Backend}{backend}} e Architettura del \textbf{PoC}.
    \item Perfezionamento dell'\textbf{AdR}.
\end{itemize}

\paragraph{\textit{Preventivo} orario ed economico}

I ruoli assegnati sono: \textbf{Responsabile}, \textbf{Amministratore}, \textbf{Analista}, \textbf{Progettista} e \textbf{Verificatore}.
\begin{table}[H]
    \centering
    \renewcommand{\arraystretch}{1.3}
    \begin{tabularx}{0.9\textwidth}{|X|c|c|c|c|c|c|c|}
        \hline
        \textbf{\textit{Sprint} 7} & \textbf{RE} & \textbf{AM} & \textbf{AN} &
        \textbf{PG} & \textbf{PR} & \textbf{VE} & \textbf{Totale} \\ \hline
        Andrea Fistarol & & & & 5 & & & \textbf{5} \\ \hline
        Dario Pirrone & & & 12 & & & & \textbf{12} \\ \hline
        Edoardo Granziero & & 5 & & & & & \textbf{5} \\ \hline
        Marco Barbiero & & & & 5 & & & \textbf{5} \\ \hline
        Samuele Bertin & 5 & & & & & & \textbf{5} \\ \hline
        Sara Ristovic & & & & & & 15 & \textbf{15} \\ \hline
        \textbf{Ore totali per ruolo} & \textbf{5} & \textbf{5} &
        \textbf{12} & \textbf{10} & \textbf{0} & \textbf{15} & \textbf{47}\\ \hline
    \end{tabularx}
    \caption{\textit{Preventivo} orario - \textit{sprint} 7}
\end{table}

Le risorse economiche previste per lo \textit{sprint} sono:

\begin{table}[H]
    \centering
    \renewcommand{\arraystretch}{1.3}
    \begin{tabularx}{0.9\textwidth}{|X|c|c|c|c|c|c|c|}
        \hline
        \textbf{\textit{Sprint} 7} & \textbf{RE} & \textbf{AM} & \textbf{AN} &
        \textbf{PG} & \textbf{PR} & \textbf{VE} & \textbf{Totale} \\ \hline
        \textbf{Costo orario} & \EUR{30} & \EUR{20} & \EUR{25} & \EUR{25} & \EUR{15}
     
        & \EUR{15} & / \\ \hline
        \textbf{Costo totale} & \EUR{150} & \EUR{100} & \EUR{300} & \EUR{250} & \EUR{0} & \EUR{225} & \textbf{\EUR{1025}} \\ \hline
    \end{tabularx}
    \caption{\textit{Preventivo} economico - \textit{sprint} 7}
\end{table}


\paragraph{\textit{Consuntivo}}
È risultato un consumo orario complessivo uguale a quanto preventivato inizialmente, ma con una ripartizione delle ore per membro diversa da quanto previsto.

\begin{table}[H]
    \centering
    \renewcommand{\arraystretch}{1.3}
    \begin{tabularx}{0.95\textwidth}{|X|c|c|c|c|c|c|c|}
        \hline
        \textbf{\textit{Sprint} 7} & \textbf{RE} & \textbf{AM} & \textbf{AN} &
        \textbf{PG} & \textbf{PR} & \textbf{VE} & \textbf{Totale} \\ \hline
        Andrea Fistarol & & & & 2 (-3) & & & \textbf{2} \\ \hline
        Dario Pirrone & & & 12 & & & & \textbf{12} \\ \hline
        Edoardo Granziero & & 7 (+2) & & & & & \textbf{7} \\ \hline
        Marco Barbiero & & & & 7 (+2) & & & \textbf{7} \\ \hline
        Samuele Bertin & 7 (+2) & & & & & & \textbf{7} \\ \hline
        Sara Ristovic & & & & & & 12 (-3) & \textbf{12} \\ \hline
        \textbf{Ore totali per ruolo} & \textbf{7} & \textbf{7} &
        \textbf{12} & \textbf{9} & \textbf{0} & \textbf{12} & \textbf{47}\\ \hline
    \end{tabularx}
    \caption{\textit{Consuntivo} orario - \textit{sprint} 7}
\end{table}

Le risorse economiche utilizzate durante lo \textit{sprint} sono:

\begin{table}[H]
    \centerline{
        \centering
        \renewcommand{\arraystretch}{1.3}
        \begin{tabularx}{1.09\textwidth}{|X|c|c|c|c|c|c|c|}
            \hline
            \textbf{\textit{Sprint} 7} & \textbf{RE} & \textbf{AM} & \textbf{AN} &
            \textbf{PG} & \textbf{PR} & \textbf{VE} & \textbf{Totale} \\ \hline
  
            \textbf{Costo orario} & \EUR{30} & \EUR{20} & \EUR{25} & \EUR{25} & \EUR{15}
            & \EUR{15} & / \\ \hline
            \textbf{Costo totale} & \EUR{210} & \EUR{140} & \EUR{300} & \EUR{225} & \EUR{0} & \EUR{180} & \textbf{\EUR{1055}} \\ \hline
            \textbf{Saldo a fine \textit{sprint}} & \textbf{\EUR{780}} & \textbf{\EUR{320}} & \textbf{\EUR{-150}} & \textbf{\EUR{2.225}} & \textbf{\EUR{2.070}} & \textbf{\EUR{1.815}} & \textbf{\EUR{7.060}}\\ \hline
        \end{tabularx}
    }
    \caption{\textit{Consuntivo} economico - \textit{sprint} 7}
\end{table}

\begin{figure}[H]
  \centering
    \includegraphics[width=0.9\textwidth]{grafici/burndown_sprint7.png}
    \caption{\textit{Grafico Burndown} - \textit{sprint} 7}
    \label{fig:burndown-sprint7}
\end{figure}

\paragraph{Analisi dei rischi occorsi}
\begin{itemize}
     \item Lo scostamento tra ore preventivate e ore effettivamente impiegate, per membro, durante questo \textit{sprint} indica che il gruppo ha riscontrato i rischi \textbf{RO4} e \textbf{RO5}, legati rispettivamente alla sottostima e alla sovrastima del tempo necessario per svolgere le proprie attività durante lo \textit{sprint}. Poiché si sono verificati scostamenti in entrambe le direzioni, il costo complessivo preventivato è stato superato di soli \EUR{30}.
     \item Nel corso di alcuni \textit{sprint} precedenti, durante le retrospettive dello \textit{sprint} concluso, il gruppo ha notato che a volte alcuni membri dichiarano di non aver completato i propri compiti, oppure concentrano il proprio impegno principalmente verso la fine dello \textit{sprint}, oppure arrivando talvolta a risultati di qualità non adeguata, tutto ciò è riconducibile al rischio \textbf{RO1}. Per aumentare la produttività di tutto il gruppo, si è deciso di provare ad accorciare la durata degli \textit{sprint}, a partire dal prossimo.
\end{itemize}

\paragraph{\textit{Retrospettiva}} Nonostante il presentarsi di diversi rischi, al termine dello \textit{sprint} tutte le \textit{issue} preventivate sono state chiuse con costi simili a quanto previsto (seppur superiori).
\paragraph{Cosa ha funzionato:}
\begin{itemize}
    \item \textbf{Buona collaborazione tra Analista e Verificatore:} durante l'intero \textit{sprint}, i membri che hanno svolto questi ruoli hanno collaborato bene e con continuità, in modalità \textit{\glo{back-and-forth}{back-and-forth}}. l'approfondimento dell'\textbf{AdR}, la verifica e l'implementazione delle correzioni avvenivano in parallelo, con una comunicazione efficace da entrambe le parti. Ciò ha contribuito in modo sostanziale alla produzione di una versione del documento di cui il gruppo è soddisfatto.
\end{itemize}

\paragraph{Migliorie per lo \textit{sprint} 8}
\begin{itemize}
    \item \textbf{Miglioramento della dichiarazione delle ore preventivate:} si deve fare una stima più precisa delle ore necessarie per ogni attività. In questo modo possiamo ridurre l'errore che si verifica tra la stima iniziale e il tempo realmente utilizzato
    \item \textbf{Accorciamento dello \textit{sprint}:} per le ragioni spiegate sopra, si è deciso di far durare il prossimo \textit{sprint} circa una settimana. Il gruppo ritiene che questa modifica possa portare diversi risultati positivi, tra cui:
    \begin{itemize}
        \item L'aumento della produttività, favorito dall'effetto psicologico di una scadenza più vicina.
        \item L'individuazione più rapida di problemi nello svolgimento di compiti, come una qualità non sufficientemente adeguata di quanto prodotto.
        \item Una comunicazione più frequente.
        \item L'individuazione più tempestiva di un impegno non sufficiente.
    \end{itemize}
    \item \textbf{Miglioramento della comunicazione interna:} per mitigare le situazioni che hanno portato alla decisione di accorciare la durata degli \textit{sprint} e per intercettarle più rapidamente qualora si ripresentassero, i membri del gruppo sono fortemente incoraggiati a comunicare in modo più frequente e tempestivo. Questo tipo di comunicazione interna comprende, tra l'altro:
    \begin{itemize}
        \item La comunicazione tempestiva di eventuali difficoltà che impediscono a un membro di completare uno o più compiti entro le scadenze stabilite
        \item La comunicazione tempestiva delle difficoltà incontrate nello svolgimento di un compito, dell'incertezza su come procedere o di dubbi sulla qualità del proprio contributo
        \item La comunicazione anticipata, e non immediatamente prima o dopo, di eventuali assenze durante lo \textit{sprint}, cercando comunque di introdurre ulteriori assenze solo quando strettamente necessario.
    \end{itemize}
\end{itemize}


\subsubsection{\textit{Sprint} 8: 08/07/2026 - 16/07/2026}

\paragraph{Obiettivi e attività}

Le attività previste per l'ottavo \textit{sprint}, ordinate in base alla priorità, sono le seguenti:

\begin{itemize}
    \item Sviluppo completo e verifica del \textbf{PoC}
    \item Perfezionamento della documentazione, ad eccezione dell'\textbf{AdR}, già completata
    \item Implementazione delle correzioni individuate durante la verifica della documentazione nello \textit{sprint} precedente.
\end{itemize}

\paragraph{\textit{Preventivo} orario ed economico}

I ruoli assegnati sono: \textbf{Responsabile}, \textbf{Amministratore}, \textbf{Programmatore} e \textbf{Verificatore}.
\begin{table}[H]
    \centering
    \renewcommand{\arraystretch}{1.3}
    \begin{tabularx}{0.9\textwidth}{|X|c|c|c|c|c|c|c|}
        \hline
        \textbf{\textit{Sprint} 8} & \multicolumn{6}{|c|}{\textbf{Ruolo}} & \\ \hline
        \textbf{Membro} & \textbf{RE} & \textbf{AM} & \textbf{AN} &
        \textbf{PG} & \textbf{PR} & \textbf{VE} & \textbf{Totale} \\ \hline
        Andrea Fistarol & & & &  & & 6 & \textbf{6} \\ \hline
        Dario Pirrone & & 4 & & & & & \textbf{4} \\ \hline
        Edoardo Granziero & & & & & & 6 & \textbf{6} \\ \hline
        Marco Barbiero & & & & & 10 & & \textbf{10} \\ \hline
        Samuele Bertin & & & & & 10 & & \textbf{10} \\ \hline
        Sara Ristovic & 5 & & & & & & \textbf{5} \\ \hline
        \textbf{Ore totali per ruolo} & \textbf{5} & \textbf{4} &
        \textbf{0} & \textbf{0} & \textbf{20} & \textbf{12} & \textbf{41}\\ \hline
    \end{tabularx}
    \caption{\textit{Preventivo} orario - \textit{sprint} 8}
\end{table}

Le risorse economiche previste per lo \textit{sprint} sono:

\begin{table}[H]
    \centering
    \renewcommand{\arraystretch}{1.3}
    \begin{tabularx}{0.9\textwidth}{|X|c|c|c|c|c|c|c|}
        \hline
        \textbf{\textit{Sprint} 8} & \textbf{RE} & \textbf{AM} & \textbf{AN} &
        \textbf{PG} & \textbf{PR} & \textbf{VE} & \textbf{Totale} \\ \hline
        \textbf{Costo orario} & \EUR{30} & \EUR{20} & \EUR{25} & \EUR{25} & \EUR{15}
     
        & \EUR{15} & / \\ \hline
        \textbf{Costo totale} & \EUR{150} & \EUR{80} & \EUR{0} & \EUR{0} & \EUR{300} & \EUR{180} & \textbf{\EUR{710}} \\ \hline
    \end{tabularx}
    \caption{\textit{Preventivo} economico - \textit{sprint} 8}
\end{table}

\paragraph{\textit{Consuntivo}}
È risultato un consumo orario maggiore di quello preventivato inizialmente
per l'\textbf{Amministratore}, mentre per i \textbf{Progettisti} e \textbf{Verificatori} il consumo complessivo risulta uguale ma con una ripartizione oraria per membro diversa da quanto preventivato.

\begin{table}[H]
    \centering
    \renewcommand{\arraystretch}{1.3}
    \begin{tabularx}{0.95\textwidth}{|X|c|c|c|c|c|c|c|}
        \hline
        \textbf{\textit{Sprint} 8} & \multicolumn{6}{|c|}{\textbf{Ruolo}} & \\ \hline
        \textbf{Membro} & \textbf{RE} & \textbf{AM} & \textbf{AN} &
        \textbf{PG} & \textbf{PR} & \textbf{VE} & \textbf{Totale} \\ \hline
        Andrea Fistarol & & & &  & & 4 (-2) & \textbf{4} \\ \hline
        Dario Pirrone & & 5 (+1) & & & & & \textbf{5} \\ \hline
        Edoardo Granziero & & & & & & 8 (+2) & \textbf{8} \\ \hline
        Marco Barbiero & & & & & 15 (+5) & & \textbf{15} \\ \hline
        Samuele Bertin & & & & & 5 (-5) & & \textbf{5} \\ \hline
        Sara Ristovic & 5 & & & & & & \textbf{5} \\ \hline
        \textbf{Ore totali per ruolo} & \textbf{5} & \textbf{5} &
        \textbf{0} & \textbf{0} & \textbf{20} & \textbf{12} & \textbf{42 (+1)}\\ \hline
    \end{tabularx}
    \caption{\textit{Consuntivo} orario - \textit{sprint} 8}
\end{table}

Le risorse economiche utilizzate durante lo \textit{sprint} sono:

\begin{table}[H]
    \centerline{
        \centering
        \renewcommand{\arraystretch}{1.3}
        \begin{tabularx}{1.09\textwidth}{|X|c|c|c|c|c|c|c|}
            \hline
            \textbf{\textit{Sprint} 8} & \textbf{RE} & \textbf{AM} & \textbf{AN} &
            \textbf{PG} & \textbf{PR} & \textbf{VE} & \textbf{Totale} \\ \hline
  
            \textbf{Costo orario} & \EUR{30} & \EUR{20} & \EUR{25} & \EUR{25} & \EUR{15}
            & \EUR{15} & / \\ \hline
            \textbf{Costo totale} & \EUR{150} & \EUR{100} & \EUR{0} & \EUR{0} & \EUR{300} & \EUR{180} & \textbf{\EUR{730}} \\ \hline
            \textbf{Saldo a fine \textit{sprint}} & \textbf{\EUR{630}} & \textbf{\EUR{220}} & \textbf{\EUR{-150}} & \textbf{\EUR{2.225}} & \textbf{\EUR{1.770}} & \textbf{\EUR{1.635}} & \textbf{\EUR{6.330}}\\ \hline
        \end{tabularx}
    }
    \caption{\textit{Consuntivo} economico - \textit{sprint} 8}
\end{table}

\begin{figure}[H]
   \centering
    \includegraphics[width=0.9\textwidth]{grafici/burndown_sprint8.png}
    \caption{\textit{Grafico Burndown} - \textit{sprint} 8}
    \label{fig:burndown-sprint8}
\end{figure}

\paragraph{Analisi dei rischi occorsi} Durante lo svolgimento delle attività di verifica in parallelo da parte di molteplici membri, il gruppo ha riscontrato il rischio di Comunicazione Interna Inefficace \textbf{(RO2)}. Nello specifico, la mancanza di una registrazione puntuale di chi stesse controllando quale sezione dello stesso documento ha causato incertezze sulla copertura dei controlli, rischiando di generare sovrapposizioni o, al contrario, di lasciare sezioni non verificate.


\paragraph{\textit{Retrospettiva}} Lo \textit{sprint} si è concluso con il completamento di tutti i compiti pianificati, stando in tempi e budget simili a quanto preventivato, mostrando la solidità dell'approccio avuto, d'altra parte si sono presentate delle nuove problematiche che che porteranno dei cambiamenti nei futuri \textit{sprint}.

\paragraph{Cosa ha funzionato:}
\begin{itemize}
    \item \textbf{\textit{Sprint} di durata 9 giorni:} Il gruppo è soddisfatto della durata ridotta dello \textit{sprint}, in quanto consente un'organizzazione più continuativa del lavoro, migliore comunicazione e maggiore efficienza.
\end{itemize}


\paragraph{Migliorie per lo \textit{sprint} 9:}
\begin{itemize}
    \item \textbf{Tracciamento formalizzato delle verifiche e responsabilità individuali:} Per snellire il lavoro in parallelo e garantire la completa tracciabilità delle revisioni, si è stabilito l'obbligo per ciascun verificatore di inserire il proprio nome nell'apposita tabella di versionamento in corrispondenza delle sezioni controllate. Questa modifica operativa permette di:
    \begin{itemize}
        \item Eliminare il rischio di sovrapposizioni e duplicazione del lavoro tra verificatori.
        \item Garantire una visibilità immediata e in tempo reale sullo stato di avanzamento delle verifiche
        \item Chiarire le responsabilità individuali, garantendo che nessuna sezione del documento venga omessa.
    \end{itemize}
\end{itemize}



\subsubsection{\textit{Sprint} 9: 17/07/2026 - 22/07/2026}

Questo \textit{Sprint} è stato interrotto la mattina del 22 luglio, anche se era stato pianificato per durare fino al 25 luglio. Ciò è successo perché il gruppo non ha superato la revisione \textbf{RTB} con il Prof. Cardin e di conseguenza ha dovuto chiudere lo \textit{Sprint} e rivedere completamente le priorità.

I seguenti obiettivi e attività erano validi fino al momento della chiusura anticipata dello \textit{Sprint}. Gli obiettivi e le attività aggiornati, stabiliti dopo la riunione \textbf{RTB} con il Prof. Cardin, sono riportati nella sezione dedicata allo \textit{Sprint} 10.

\paragraph{Obiettivi e attività}

Le attività previste per il nono \textit{Sprint}, ordinate in base alla priorità, sono le seguenti:

\begin{itemize}
    \item Preparazione alla revisione \textbf{RTB} con il Prof. Cardin: creazione delle slide ed esercitazione dell'esposizione dei vari contenuti
    \item Controllo dei collegamenti tra documenti e dei termini del \textbf{Glossario}, oltre a un controllo contenutistico di alcune sezioni
    \item Aggiornamento del sito e dei vari PDF dei documenti esposti su di esso, portando il documento di \textbf{AdR} alla versione 1.0.0.
\end{itemize}

\paragraph{\textit{Preventivo} orario ed economico}

I ruoli assegnati sono: \textbf{Responsabile}, \textbf{Amministratore}, \textbf{Programmatore} e \textbf{Verificatore}.
\begin{table}[H]
    \centering
    \renewcommand{\arraystretch}{1.3}
    \begin{tabularx}{0.9\textwidth}{|X|c|c|c|c|c|c|c|}
        \hline
        \textbf{\textit{Sprint} 9} & \multicolumn{6}{|c|}{\textbf{Ruolo}} & \\ \hline
        \textbf{Membro} & \textbf{RE} & \textbf{AM} & \textbf{AN} & \textbf{PG} & \textbf{PR} & \textbf{VE} & \textbf{Totale} \\ \hline
        Andrea Fistarol  & & & & & & 4 & \textbf{4} \\ \hline
        Dario Pirrone    & & & & & & 7 & \textbf{7} \\ \hline
        Edoardo Granziero & & & & & 5 & & \textbf{5} \\ \hline
        Marco Barbiero   & 2 & & & & & & \textbf{2} \\ \hline
        Samuele Bertin   & & & & & & 6 & \textbf{6} \\ \hline
        Sara Ristovic    & & 4 & & & & & \textbf{4} \\ \hline
        \textbf{Ore totali per ruolo} & \textbf{2} & \textbf{4} & \textbf{0} & \textbf{0} & \textbf{5} & \textbf{17} & \textbf{28} \\ \hline
    \end{tabularx}
    \caption{\textit{Preventivo} orario - \textit{Sprint} 9}
\end{table}

Le risorse economiche previste per lo \textit{sprint} sono:

\begin{table}[H]
    \centering
    \renewcommand{\arraystretch}{1.3}
    \begin{tabularx}{0.9\textwidth}{|X|c|c|c|c|c|c|c|}
        \hline
        \textbf{\textit{Sprint} 9} & \textbf{RE} & \textbf{AM} & \textbf{AN} &
        \textbf{PG} & \textbf{PR} & \textbf{VE} & \textbf{Totale} \\ \hline
        \textbf{Costo orario} & \EUR{30} & \EUR{20} & \EUR{25} & \EUR{25} & \EUR{15}
     
        & \EUR{15} & / \\ \hline
        \textbf{Costo totale} & \EUR{60} & \EUR{80} & \EUR{0} & \EUR{0} & \EUR{75} & \EUR{255} & \textbf{\EUR{470}} \\ \hline
    \end{tabularx}
    \caption{\textit{Preventivo} economico - \textit{sprint} 9}
\end{table}

\paragraph{\textit{Consuntivo}}
È risultato un consumo orario complessivo significativamente inferiore a quanto preventivato inizialmente, a causa della chiusura anticipata dello \textit{Sprint}.

\begin{table}[H]
    \centering
    \renewcommand{\arraystretch}{1.3}
    \begin{tabularx}{0.95\textwidth}{|X|c|c|c|c|c|c|c|}
        \hline
        \textbf{\textit{Sprint} 9} & \multicolumn{6}{|c|}{\textbf{Ruolo}} & \\ \hline
        \textbf{Membro} & \textbf{RE} & \textbf{AM} & \textbf{AN} &
        \textbf{PG} & \textbf{PR} & \textbf{VE} & \textbf{Totale} \\ \hline
        Andrea Fistarol & & & & & & 4 & \textbf{4} \\ \hline
        Dario Pirrone & & & & & & 3 (-4) & \textbf{3} \\ \hline
        Edoardo Granziero & & & & & 1 (-4) & & \textbf{1} \\ \hline
        Marco Barbiero & 3 (+1) & & & & & & \textbf{3} \\ \hline
        Samuele Bertin & & & & & & 2 (-4) & \textbf{2} \\ \hline
        Sara Ristovic & & 2 (-2) & & & & & \textbf{2} \\ \hline
        \textbf{Ore totali per ruolo} & \textbf{3} & \textbf{2} &
        \textbf{0} & \textbf{0} & \textbf{1} & \textbf{9} & \textbf{15 (-13)}\\ \hline
    \end{tabularx}
    \caption{\textit{Consuntivo} orario - \textit{sprint} 9}
\end{table}

Le risorse economiche utilizzate durante lo \textit{sprint} sono:

\begin{table}[H]
    \centerline{
        \centering
        \renewcommand{\arraystretch}{1.3}
        \begin{tabularx}{1.09\textwidth}{|X|c|c|c|c|c|c|c|}
            \hline
            \textbf{\textit{Sprint} 9} & \textbf{RE} & \textbf{AM} & \textbf{AN} &
            \textbf{PG} & \textbf{PR} & \textbf{VE} & \textbf{Totale} \\ \hline
  
            \textbf{Costo orario} & \EUR{30} & \EUR{20} & \EUR{25} & \EUR{25} & \EUR{15}
            & \EUR{15} & / \\ \hline
            \textbf{Costo totale} & \EUR{90} & \EUR{40} & \EUR{0} & \EUR{0} & \EUR{15} & \EUR{135} & \textbf{\EUR{280}} \\ \hline
            \textbf{Saldo a fine \textit{sprint}} & \textbf{\EUR{540}} & \textbf{\EUR{180}} & \textbf{\EUR{-150}} & \textbf{\EUR{2.225}} & \textbf{\EUR{1.755}} & \textbf{\EUR{1.500}} & \textbf{\EUR{6.050}}\\ \hline
        \end{tabularx}
    }
    \caption{\textit{Consuntivo} economico - \textit{sprint} 9}
\end{table}

\begin{figure}[H]
  \centering
    \includegraphics[width=0.9\textwidth]{grafici/burndown_sprint9.png}
    \caption{\textit{Grafico Burndown} - \textit{sprint} 9}
    \label{fig:burndown-sprint9}
\end{figure}

\paragraph{Analisi dei rischi occorsi}
\begin{itemize}
     \item \textbf{Semaforo rosso all'\textbf{RTB}:} Il gruppo non ha superato la revisione \textbf{RTB} con il Prof. Cardin e, di conseguenza, ha dovuto fermarsi e procedere con il miglioramento del \textbf{PoC}. Questo ha prodotto un ritardo inatteso nell'avanzamento del progetto.
\end{itemize}

\paragraph{\textit{Retrospettiva}} Poiché lo \textit{Sprint} è stato chiuso in anticipo, molti compiti non sono stati completati e sono stati temporaneamente sospesi o rinviati. Saranno completati quando i problemi individuati nel \textbf{PoC} saranno stati risolti.

\paragraph{Migliorie per lo \textit{Sprint} 10}
\begin{itemize}
    \item \textbf{Comunicazione anticipata con i professori prima delle revisioni:}  Come suggerito anche dal Prof. Vardanega, per mitigare il rischio che si verifichi nuovamente un semaforo rosso in futuro, il gruppo è incoraggiato a contattare i professori e chiedere domande e riscontro prima delle revisioni stesse. Ciò aiuterebbe il gruppo a identificare i problemi critici prima delle revisioni, evitando che possano portare a un semaforo rosso oppure a un voto basso.
\end{itemize}


\subsubsection{\textit{Sprint} 10: 23/07/2026 - 04/08/2026}

\paragraph{Obiettivi e attività}

Le attività previste per il decimo \textit{sprint}, ordinate in base alla priorità, sono le seguenti:

\begin{itemize}
    \item riprogettazione del \textbf{PoC}, in modo da incorporare il maggior numero possibile delle tecnologie scelte
    \item codifica delle componenti nuove del \textbf{PoC}
\end{itemize}

\paragraph{\textit{Preventivo} orario ed economico}

I ruoli assegnati sono: \textbf{Responsabile}, \textbf{Amministratore}, \textbf{Progettista}, \textbf{Programmatore} e \textbf{Verificatore}. In questo \textit{sprint} si sperimenta per la prima volta l'assegnazione di più ruoli allo stesso membro del gruppo.
\begin{table}[H]
    \centering
    \renewcommand{\arraystretch}{1.3}
    \begin{tabularx}{0.9\textwidth}{|X|c|c|c|c|c|c|c|}
        \hline
        \textbf{\textit{Sprint} 10} & \multicolumn{6}{|c|}{\textbf{Ruolo}} & \\ \hline
        \textbf{Membro} & \textbf{RE} & \textbf{AM} & \textbf{AN} &
        \textbf{PG} & \textbf{PR} & \textbf{VE} & \textbf{Totale} \\ \hline
        Andrea Fistarol & & & & & 9 & & \textbf{9} \\ \hline
        Dario Pirrone & & & & & 13 & & \textbf{13} \\ \hline
        Edoardo Granziero & & 1 & & & & 3 & \textbf{4} \\ \hline
        Marco Barbiero & & & & 8 & & & \textbf{8} \\ \hline
        Samuele Bertin & & & & 10 & & & \textbf{10} \\ \hline
        Sara Ristovic & 1 & & & & & 5 & \textbf{6} \\ \hline
        \textbf{Ore totali per ruolo} & \textbf{1} & \textbf{1} &
        \textbf{0} & \textbf{18} & \textbf{22} & \textbf{8} & \textbf{50}\\ \hline
    \end{tabularx}
    \caption{\textit{Preventivo} orario - \textit{sprint} 10}
\end{table}

Le risorse economiche previste per lo \textit{sprint} sono:

\begin{table}[H]
    \centering
    \renewcommand{\arraystretch}{1.3}
    \begin{tabularx}{0.9\textwidth}{|X|c|c|c|c|c|c|c|}
        \hline
        \textbf{\textit{Sprint} 10} & \textbf{RE} & \textbf{AM} & \textbf{AN} &
        \textbf{PG} & \textbf{PR} & \textbf{VE} & \textbf{Totale} \\ \hline
        \textbf{Costo orario} & \EUR{30} & \EUR{20} & \EUR{25} & \EUR{25} & \EUR{15}
     
        & \EUR{15} & / \\ \hline
        \textbf{Costo totale} & \EUR{30} & \EUR{20} & \EUR{0} & \EUR{450} & \EUR{330} & \EUR{120} & \textbf{\EUR{950}} \\ \hline
    \end{tabularx}
    \caption{\textit{Preventivo} economico - \textit{sprint} 10}
\end{table}


\paragraph{\textit{Consuntivo}}
È risultato un consumo orario complessivo lievemente inferiore a quanto preventivato, in generale con meno ore spese nel ruolo di \textbf{Verificatore} e lievemente di più per il ruolo di \textbf{Programmatore}.

\begin{table}[H]
    \centering
    \renewcommand{\arraystretch}{1.3}
    \begin{tabularx}{0.95\textwidth}{|X|c|c|c|c|c|c|c|}
        \hline
        \textbf{\textit{Sprint} 10} & \multicolumn{6}{|c|}{\textbf{Ruolo}} & \\ \hline
        \textbf{Membro} & \textbf{RE} & \textbf{AM} & \textbf{AN} &
        \textbf{PG} & \textbf{PR} & \textbf{VE} & \textbf{Totale} \\ \hline
        Andrea Fistarol & & & & & 8 (-1) & & \textbf{8} \\ \hline
        Dario Pirrone & & & & & 16 (+3) & & \textbf{16} \\ \hline
        Edoardo Granziero & & 0 (-1) & & & & 2 (-1) & \textbf{2} \\ \hline
        Marco Barbiero & & & & 8 & & & \textbf{8} \\ \hline
        Samuele Bertin & & & & 10 & & & \textbf{10} \\ \hline
        Sara Ristovic & 1 & & & & & 3 (-2) & \textbf{4} \\ \hline
        \textbf{Ore totali per ruolo} & \textbf{1} & \textbf{0} &
        \textbf{0} & \textbf{18} & \textbf{24} & \textbf{5} & \textbf{48 (-2)}\\ \hline
    \end{tabularx}
    \caption{\textit{Consuntivo} orario - \textit{sprint} 10}
\end{table}

Le risorse economiche utilizzate durante lo \textit{sprint} sono:

\begin{table}[H]
    \centerline{
        \centering
        \renewcommand{\arraystretch}{1.3}
        \begin{tabularx}{1.09\textwidth}{|X|c|c|c|c|c|c|c|}
            \hline
            \textbf{\textit{Sprint} 10} & \textbf{RE} & \textbf{AM} & \textbf{AN} &
            \textbf{PG} & \textbf{PR} & \textbf{VE} & \textbf{Totale} \\ \hline
  
            \textbf{Costo orario} & \EUR{30} & \EUR{20} & \EUR{25} & \EUR{25} & \EUR{15}
            & \EUR{15} & / \\ \hline
            \textbf{Costo totale} & \EUR{30} & \EUR{0} & \EUR{0} & \EUR{450} & \EUR{360} & \EUR{75} & \textbf{\EUR{915}} \\ \hline
            \textbf{Saldo a fine \textit{sprint}} & \textbf{\EUR{510}} & \textbf{\EUR{180}} & \textbf{\EUR{-150}} & \textbf{\EUR{1.775}} & \textbf{\EUR{1.395}} & \textbf{\EUR{1.425}} & \textbf{\EUR{5.135}}\\ \hline
        \end{tabularx}
    }
    \caption{\textit{Consuntivo} economico - \textit{sprint} 10}
\end{table}

\begin{figure}[H]
  \centering
    \includegraphics[width=0.9\textwidth]{grafici/burndown_sprint10.png}
    \caption{\textit{Grafico Burndown} - \textit{sprint} 10}
    \label{fig:burndown-sprint10}
\end{figure}

\paragraph{Analisi dei rischi occorsi}
\begin{itemize}
    \item \textbf{Scarsa Pianificazione (RO1) e Comunicazione Interna Inefficace (RO2):} Durante lo svolgimento dello \textit{sprint}, l'organizzazione iniziale ha visto impegnati a pieno regime solo i membri assegnati ai ruoli di progettazione e poi di codifica, lasciando il resto del gruppo con compiti minimi da svolgere in quei giorni. Questa mancanza di pianificazione trasversale e di comunicazione ha portato all'accumulo di ritardi sul lungo termine, seppure minimi.
\end{itemize}

\paragraph{\textit{Retrospettiva}} Lo \textit{sprint} si è concluso positivamente per quanto riguarda le tempistiche, chiudendo nei tempi previsti tutte le \textit{issue}. L'obiettivo principale è stato raggiunto con l'ultimazione del nuovo \textbf{Proof of Concept} (\textbf{PoC}), sebbene si siano evidenziate nuovamente delle lacune organizzative interne che andranno corrette immediatamente.

\paragraph{Cosa ha funzionato:}
\begin{itemize}
    \item \textbf{Solidità del nuovo PoC:} Nonostante le difficoltà organizzative, il team è  soddisfatto del \textbf{PoC} realizzato. Il lavoro svolto non andrà perso, ma rappresenterà una solida base di partenza per lo sviluppo del \textbf{Minimum Viable Product} (\textbf{MVP}), permettendo di riciclare diverso codice senza la necessità di ripartire da capo.
\end{itemize}

\paragraph{Migliorie per lo \textit{Sprint} 11}
\begin{itemize}
    \item \textbf{Maggiore comunicazione e flessibilità:} È fondamentale incrementare la comunicazione all'interno del gruppo e la flessibilità a cambiamenti improvvisi. Per evitare il ripresentarsi dei rischi (\textbf{RO1} e \textbf{RO2}), si è deciso di stabilirsi su \textit{sprint} di durata breve e flessibile alle necessità del momento; questo approccio aiuta a mantenere un contatto più costante tra i membri e a ripartire il carico di lavoro in modo più dinamico.
\end{itemize}


\subsubsection{\textit{Sprint} 11: 05/08/2026 - 07/08/2026}

\paragraph{Obiettivi e attività}

Questo \textit{sprint} sarà particolarmente corto e verticalmente incentrato sulla chiusura dei lavori in vista della revisione \textbf{RTB}. Le attività previste, ordinate in base alla priorità, sono le seguenti:

\begin{itemize}
    \item Preparazione e svolgimento dell'incontro di dimostrazione del \textbf{PoC} con il Prof. Cardin
    \item Ripresa e completamento delle \textit{issue} di verifica e correzione lasciate nel \textit{backlog} dallo \textit{sprint} 9, interrotto prematuramente (su \textbf{Glossario}, \textbf{NdP}, \textbf{PdQ} e \textbf{PdP})
    \item Aggiornamento del sito e stesura della lettera di candidatura per la revisione \textbf{RTB} con il Prof. Tullio
\end{itemize}

\paragraph{\textit{Preventivo} orario ed economico}

I ruoli assegnati per questo breve \textit{sprint} conclusivo sono: \textbf{Responsabile}, \textbf{Amministratore} e \textbf{Verificatore}. Essendo uno \textit{sprint} di rifinitura e correzione documenti, non sono previste ore per i ruoli di Analista, Progettista e Programmatore.

\begin{table}[H]
    \centering
    \renewcommand{\arraystretch}{1.3}
    \begin{tabularx}{0.9\textwidth}{|X|c|c|c|c|c|c|c|}
        \hline
        \textbf{\textit{Sprint} 11} & \multicolumn{6}{|c|}{\textbf{Ruolo}} & \\ \hline
        \textbf{Membro} & \textbf{RE} & \textbf{AM} & \textbf{AN} &
        \textbf{PG} & \textbf{PR} & \textbf{VE} & \textbf{Totale} \\ \hline
        Andrea Fistarol & & & & & & 6 & \textbf{6} \\ \hline
        Dario Pirrone & 2 & & & & & & \textbf{2} \\ \hline
        Edoardo Granziero & & & & & & 2 & \textbf{2} \\ \hline
        Marco Barbiero & & 3 & & & & & \textbf{3} \\ \hline
        Samuele Bertin & & & & & & 6 & \textbf{6} \\ \hline
        Sara Ristovic & & & & & & 3 & \textbf{3} \\ \hline
        \textbf{Ore totali per ruolo} & \textbf{2} & \textbf{3} & \textbf{0} & \textbf{0} & \textbf{0} & \textbf{17} & \textbf{22}\\ \hline
    \end{tabularx}
    \caption{\textit{Preventivo} orario - \textit{sprint} 11}
\end{table}

Le risorse economiche previste per lo \textit{sprint} sono:

\begin{table}[H]
    \centering
    \renewcommand{\arraystretch}{1.3}
    \begin{tabularx}{0.9\textwidth}{|X|c|c|c|c|c|c|c|}
        \hline
        \textbf{\textit{Sprint} 11} & \textbf{RE} & \textbf{AM} & \textbf{AN} &
        \textbf{PG} & \textbf{PR} & \textbf{VE} & \textbf{Totale} \\ \hline
        \textbf{Costo orario} & \EUR{30} & \EUR{20} & \EUR{25} & \EUR{25} & \EUR{15} & \EUR{15} & / \\ \hline
        \textbf{Costo totale} & \EUR{60} & \EUR{60} & \EUR{0} & \EUR{0} & \EUR{0} & \EUR{255} & \textbf{\EUR{375}} \\ \hline
    \end{tabularx}
    \caption{\textit{Preventivo} economico - \textit{sprint} 11}
\end{table}

\paragraph{\textit{Consuntivo}}
È risultato un consumo orario complessivo quasi corrispondente al preventivo, con solo un'ora in più per il ruolo di \textbf{Verificatore}.

\begin{table}[H]
    \centering
    \renewcommand{\arraystretch}{1.3}
    \begin{tabularx}{0.95\textwidth}{|X|c|c|c|c|c|c|c|}
        \hline
        \textbf{\textit{Sprint} 11} & \multicolumn{6}{|c|}{\textbf{Ruolo}} & \\ \hline
        \textbf{Membro} & \textbf{RE} & \textbf{AM} & \textbf{AN} &
        \textbf{PG} & \textbf{PR} & \textbf{VE} & \textbf{Totale} \\ \hline
        Andrea Fistarol & & & & & & 5 (-1) & \textbf{5} \\ \hline
        Dario Pirrone & 2 & & & & & & \textbf{2} \\ \hline
        Edoardo Granziero & & & & & & 4 (+2) & \textbf{4} \\ \hline
        Marco Barbiero & & 3 & & & & & \textbf{3} \\ \hline
        Samuele Bertin & & & & & & 6 & \textbf{6} \\ \hline
        Sara Ristovic & 1 & & & & & 3 & \textbf{3} \\ \hline
        \textbf{Ore totali per ruolo} & \textbf{2} & \textbf{3} &
        \textbf{0} & \textbf{0} & \textbf{0} & \textbf{18 (+1)} & \textbf{23 (+1)}\\ \hline
    \end{tabularx}
    \caption{\textit{Consuntivo} orario - \textit{sprint} 11}
\end{table}

Le risorse economiche utilizzate durante lo \textit{sprint} sono:

\begin{table}[H]
    \centerline{
        \centering
        \renewcommand{\arraystretch}{1.3}
        \begin{tabularx}{1.09\textwidth}{|X|c|c|c|c|c|c|c|}
            \hline
            \textbf{\textit{Sprint} 11} & \textbf{RE} & \textbf{AM} & \textbf{AN} &
            \textbf{PG} & \textbf{PR} & \textbf{VE} & \textbf{Totale} \\ \hline
  
            \textbf{Costo orario} & \EUR{30} & \EUR{20} & \EUR{25} & \EUR{25} & \EUR{15}
            & \EUR{15} & / \\ \hline
            \textbf{Costo totale} & \EUR{60} & \EUR{60} & \EUR{0} & \EUR{0} & \EUR{0} & \EUR{270} & \textbf{\EUR{390}} \\ \hline
            \textbf{Saldo a fine \textit{sprint}} & \textbf{\EUR{450}} & \textbf{\EUR{120}} & \textbf{\EUR{-150}} & \textbf{\EUR{1.775}} & \textbf{\EUR{1.395}} & \textbf{\EUR{1.115}} & \textbf{\EUR{4.745}}\\ \hline
        \end{tabularx}
    }
    \caption{\textit{Consuntivo} economico - \textit{sprint} 11}
\end{table}

\begin{figure}[H]
  \centering
    \includegraphics[width=0.9\textwidth]{grafici/burndown_sprint11.png}
    \caption{\textit{Grafico Burndown} - \textit{sprint} 11}
    \label{fig:burndown-sprint11}
\end{figure}

\paragraph{\textit{Retrospettiva}} Lo \textit{sprint} è stato molto corto (4 gg) ma produttivo e, come dimostra il grafico \textit{Burndown}, con una buona gestione dei compiti nel tempo.

\paragraph{Cosa ha funzionato:}
\begin{itemize}
    \item \textbf{Gestione dei compiti:} le \textit{issue} avevano adeguata capillarità e sono state affrontate senza ritardi, cosa probabilmente dovuta anche alla scadenza ravvicinata.
\end{itemize}

\paragraph{Migliorie per lo \textit{Sprint} 12}
\begin{itemize}
    \item \textbf{Mantenere alto il livello di impegno:} lo \textit{sprint} 12 sarà il primo della fase \textbf{PB}. Con cognizione del ritardo accumulato, sarà importante far valere al meglio il tempo a disposizione, nell'ottica di non introdurre ulteriore ritardo ma anzi di riuscire a recuperarne una parte.
\end{itemize}

\newpage
\subsubsection{\textit{Consuntivo} \textbf{RTB}}
\paragraph{Ore spese} ~\par
\begin{table}[H]
    \centering
    \renewcommand{\arraystretch}{1.3}
    \begin{tabularx}{0.95\textwidth}{|X|c|c|c|c|c|c|c|}
        \hline
        \textbf{RTB} & \multicolumn{6}{|c|}{\textbf{Ruolo}} & \\ \hline
        \textbf{Membro} & \textbf{RE} & \textbf{AM} & \textbf{AN} & \textbf{PG} & \textbf{PR} & \textbf{VE} & \textbf{Totale} \\ \hline
        Andrea Fistarol & 1 & 5 & 11 & 3 & 8 & 13 & \textbf{41} \\ \hline
        Dario Pirrone & 6 & 9 & 26 & 0 & 16 & 6 & \textbf{63} \\ \hline
        Edoardo Granziero & 4 & 12 & 3 & 0 & 1 & 19 & \textbf{39} \\ \hline
        Marco Barbiero & 3 & 6 & 8 & 24 & 15 & 4 & \textbf{60} \\ \hline
        Samuele Bertin & 14 & 4 & 4 & 10 & 5 & 12 & \textbf{49} \\ \hline
        Sara Ristovic & 11 & 6 & 20 & 0 & 0 & 19 & \textbf{56} \\ \hline
        \textbf{Ore totali per ruolo} & \textbf{39} & \textbf{42} & \textbf{72} & \textbf{37} & \textbf{45} & \textbf{73} & \textbf{308}\\ \hline
    \end{tabularx}
    \caption{\textit{Consuntivo} ore spese - \textbf{RTB}}
\end{table}

\begin{figure}[H]
  \centering
    \includegraphics[width=0.9\textwidth]{grafici/torta_membri.png}
    \caption{Suddivisione oraria per membro - \textbf{RTB}}
    \label{fig:membriRTB}
\end{figure}

\begin{figure}[H]
  \centering
    \includegraphics[width=0.9\textwidth]{grafici/torta_ruoli.png}
    \caption{Suddivisione oraria per ruolo - \textbf{RTB}}
    \label{fig:ruoliRTB}
\end{figure}

\paragraph{Ore rimanenti} ~\par
\begin{table}[H]
    \centering
    \renewcommand{\arraystretch}{1.3}
    \begin{tabularx}{0.95\textwidth}{|X|c|c|c|c|c|c|c|}
        \hline
        \textbf{RTB} & \multicolumn{6}{|c|}{\textbf{Ruolo}} & \\ \hline
        \textbf{Membro} & \textbf{RE} & \textbf{AM} & \textbf{AN} & \textbf{PG} & \textbf{PR} & \textbf{VE} & \textbf{Totale} \\ \hline
        Andrea Fistarol & 8 & 3 & 0 & 15 & 15 & 12 & \textbf{53} \\ \hline
        Dario Pirrone & 3 & -1 & -15 & 18 & 7 & 19 & \textbf{31} \\ \hline
        Edoardo Granziero & 5 & -4 & 8 & 18 & 22 & 6 & \textbf{55} \\ \hline
        Marco Barbiero & 6 & 2 & 3 & -6 & 8 & 21 & \textbf{34} \\ \hline
        Samuele Bertin & -5 & 4 & 7 & 8 & 18 & 13 & \textbf{45} \\ \hline
        Sara Ristovic & -2 & 2 & -9 & 18 & 23 & 6 & \textbf{38} \\ \hline
        \textbf{Ore totali per ruolo} & \textbf{15} & \textbf{6} & \textbf{-6} & \textbf{71} & \textbf{93} & \textbf{77} & \textbf{256}\\ \hline
    \end{tabularx}
    \caption{\textit{Consuntivo} ore rimanenti - \textbf{RTB}}
\end{table}

\paragraph{Costi affrontati} ~\par
\begin{table}[H]
    \centerline{
        \centering
        \renewcommand{\arraystretch}{1.3}
        \begin{tabularx}{1.09\textwidth}{|X|c|c|c|c|c|c|c|}
            \hline
            \textbf{RTB} & \textbf{RE} & \textbf{AM} & \textbf{AN} & \textbf{PG} & \textbf{PR} & \textbf{VE} & \textbf{Totale} \\ \hline  
            \textbf{Costo orario} & \EUR{30} & \EUR{20} & \EUR{25} & \EUR{25} & \EUR{15} & \EUR{15} & / \\ \hline
            \textbf{Costo totale} & \EUR{1.170} & \EUR{840} & \EUR{1.800} & \EUR{925} & \EUR{975} & \EUR{1.095} & \textbf{\EUR{6.505}} \\ \hline
            \textbf{Saldo a fine \textbf{RTB}} & \textbf{\EUR{450}} & \textbf{\EUR{120}} & \textbf{\EUR{-150}} & \textbf{\EUR{1.775}} & \textbf{\EUR{1.395}} & \textbf{\EUR{1.115}} & \textbf{\EUR{4.745}}\\ \hline
        \end{tabularx}
    }
    \caption{\textit{Consuntivo} economico - \textbf{RTB}}
\end{table}

\paragraph{Considerazioni} I dati del \textit{consuntivo}, uniti a quelli del \href{https://skynetunigroup.github.io/Code_Guardian/Documentazione/RTB/piano_di_qualifica_v1.0.0.pdf}{\textit{Cruscotto di Verifica}}, indicano un eccesso di costi rispetto al valore raggiunto e un ritardo pericoloso. La situazione, però, non è irrecuperabile. Per il problema del budget, verrà attuata una riduzione dello scope del progetto con l'esclusione di alcuni requisiti non obbligatori dall'\textbf{MVP}; un fattore positivo sarà la preponderanza, nel prossimo periodo, dei ruoli del \textbf{Programmatore} e del \textbf{Verificatore} con il loro costo orario ridotto (\EUR{15}). Per un ulteriore contenimento dei costi e per evitare di accumulare altro ritardo, o persino per recuperarne una parte, è necessario un aumento della produttività oraria, affiancato da adeguata pianificazione e organizzazione.

\subsection{Product Baseline (\textbf{PB})}

\subsubsection{\textit{Sprint} 12: 08/08/2026 - 14/08/2026}

\paragraph{Obiettivi e attività}

Questo \textit{sprint} avrà una lunghezza standard (7 gg), dopo l'eccezione di quello precedente. Le attività saranno incentrate sull'avvio della fase \textbf{PB}, con la progettazione del sistema, e sull'incontro col prof. Vardanega. Di seguito le attività individuate:

\begin{itemize}
    \item Ultimo controllo formale di correttezza della documentazione; preparazione e richiesta di incontro con il prof. Vardanega.
    \item Decidere quali requisiti desiderabili o opzionali soddisfare.
    \item Progettazione del sistema, partendo da valutazioni sul riuso di parti o soluzioni del \textbf{PoC} e integrando le parti escluse, come l'\textbf{orchestratore} e l'infrastruttura \textbf{AWS}.
\end{itemize}

\paragraph{\textit{Preventivo} orario ed economico}

I ruoli assegnati sono: \textbf{Responsabile}, \textbf{Amministratore}, \textbf{Progettista} e \textbf{Verificatore}. Il focus dello \textit{sprint} è sulla progettazione, questo motiva il gran numero di ore assegnate ai \textbf{Progettisti} e l'assenza di \textbf{Programmatori}.

\begin{table}[H]
    \centering
    \renewcommand{\arraystretch}{1.3}
    \begin{tabularx}{0.9\textwidth}{|X|c|c|c|c|c|c|c|}
        \hline
        \textbf{\textit{Sprint} 12} & \multicolumn{6}{|c|}{\textbf{Ruolo}} & \\ \hline
        \textbf{Membro} & \textbf{RE} & \textbf{AM} & \textbf{AN} & \textbf{PG} & \textbf{PR} & \textbf{VE} & \textbf{Totale} \\ \hline
        Andrea Fistarol & 3 & & & & &  & \textbf{3} \\ \hline
        Dario Pirrone & & & & 10 & & & \textbf{10} \\ \hline
        Edoardo Granziero & & & & & & 4 & \textbf{4} \\ \hline
        Marco Barbiero & & 2 & & & & 7 & \textbf{9} \\ \hline
        Samuele Bertin & & & & 10 & & & \textbf{10} \\ \hline
        Sara Ristovic & & & & 10 & & & \textbf{10} \\ \hline
        \textbf{Ore totali per ruolo} & \textbf{3} & \textbf{2} & \textbf{0} & \textbf{30} & \textbf{0} & \textbf{11} & \textbf{46}\\ \hline
    \end{tabularx}
    \caption{\textit{Preventivo} orario - \textit{sprint} 12}
\end{table}

Le risorse economiche previste per lo \textit{sprint} sono:

\begin{table}[H]
    \centering
    \renewcommand{\arraystretch}{1.3}
    \begin{tabularx}{0.9\textwidth}{|X|c|c|c|c|c|c|c|}
        \hline
        \textbf{\textit{Sprint} 12} & \textbf{RE} & \textbf{AM} & \textbf{AN} &
        \textbf{PG} & \textbf{PR} & \textbf{VE} & \textbf{Totale} \\ \hline
        \textbf{Costo orario} & \EUR{30} & \EUR{20} & \EUR{25} & \EUR{25} & \EUR{15} & \EUR{15} & / \\ \hline
        \textbf{Costo totale} & \EUR{90} & \EUR{40} & \EUR{0} & \EUR{750} & \EUR{0} & \EUR{165} & \textbf{\EUR{1.045}} \\ \hline
    \end{tabularx}
    \caption{\textit{Preventivo} economico - \textit{sprint} 12}
\end{table}

\paragraph{\textit{Consuntivo}}
È stata registrata una notevole carenza di ore rispetto al preventivo, ben 11 in meno. In parte sono dovute a una sovrastima dei tempi necessari ma principalmente sono da imputare a un progresso più difficoltoso del previsto da parte di Ristovic.

\begin{table}[H]
    \centering
    \renewcommand{\arraystretch}{1.3}
    \begin{tabularx}{0.95\textwidth}{|X|c|c|c|c|c|c|c|}
        \hline
        \textbf{\textit{Sprint} 12} & \multicolumn{6}{|c|}{\textbf{Ruolo}} & \\ \hline
        \textbf{Membro} & \textbf{RE} & \textbf{AM} & \textbf{AN} & \textbf{PG} & \textbf{PR} & \textbf{VE} & \textbf{Totale} \\ \hline
        Andrea Fistarol & 3 & & & & & & \textbf{3} \\ \hline
        Dario Pirrone & & & & 9 (-1) & & & \textbf{9} \\ \hline
        Edoardo Granziero & & & & & & 4 & \textbf{4} \\ \hline
        Marco Barbiero & & 1 (-1) & & & & 5 (-2) & \textbf{6} \\ \hline
        Samuele Bertin & & & & 10 & & & \textbf{10} \\ \hline
        Sara Ristovic & & & & 3 (-7) & & & \textbf{3} \\ \hline
        \textbf{Ore totali per ruolo} & \textbf{3} & \textbf{1} & \textbf{0} & \textbf{22} & \textbf{0} & \textbf{9} & \textbf{35 (-11)}\\ \hline
    \end{tabularx}
    \caption{\textit{Consuntivo} orario - \textit{sprint} 12}
\end{table}

Le risorse economiche utilizzate durante lo \textit{sprint} sono:

\begin{table}[H]
    \centerline{
        \centering
        \renewcommand{\arraystretch}{1.3}
        \begin{tabularx}{1.09\textwidth}{|X|c|c|c|c|c|c|c|}
            \hline
            \textbf{\textit{Sprint} 12} & \textbf{RE} & \textbf{AM} & \textbf{AN} & \textbf{PG} & \textbf{PR} & \textbf{VE} & \textbf{Totale} \\ \hline  
            \textbf{Costo orario} & \EUR{30} & \EUR{20} & \EUR{25} & \EUR{25} & \EUR{15} & \EUR{15} & / \\ \hline
            \textbf{Costo totale} & \EUR{90} & \EUR{20} & \EUR{0} & \EUR{550} & \EUR{0} & \EUR{135} & \textbf{\EUR{795}} \\ \hline
            \textbf{Saldo a fine \textit{sprint}} & \textbf{\EUR{360}} & \textbf{\EUR{100}} & \textbf{\EUR{-150}} & \textbf{\EUR{1.225}} & \textbf{\EUR{1.395}} & \textbf{\EUR{1.020}} & \textbf{\EUR{3.950}}\\ \hline
        \end{tabularx}
    }
    \caption{\textit{Consuntivo} economico - \textit{sprint} 12}
\end{table}

\begin{figure}[H]
  \centering
    \includegraphics[width=0.9\textwidth]{grafici/burndown_sprint12.png}
    \caption{\textit{Grafico Burndown} - \textit{sprint} 12}
    \label{fig:burndown-sprint12}
\end{figure}

\paragraph{Analisi dei rischi occorsi}
\begin{itemize}
    \item \textbf{Comunicazione Interna Inefficace (R02):} alcune difficoltà sono emerse tardivamente.
    \item \textbf{Sottostima delle Attività (R04):} alcune attività hanno richiesto più tempo del previsto.
\end{itemize}

\paragraph{\textit{Retrospettiva}} Il grafico \textit{Burndown} mostra un procedimento non ideale dello \textit{sprint}, molti compiti sono stati conclusi solo l'ultimo giorno, togliendo spazio alla verifica. È stato uno \textit{sprint} ambizioso, in cui è stato previsto di impiegare tante ore per la progettazione per cercare di rispettare la scadenza del 19 agosto. Da una parte, ne sarebbero servite ancora di più per completare anche l'unico compito davvero avanzato, dall'altra è mancata produttività. 3 delle 4 \textit{issue} non completate dipendevano da condizioni che non si sono verificate nel corso dello \textit{sprint} ovvero la ricezione formale del "semaforo verde" dopo il coloquio \textbf{TB} e il successivo colloquio \textbf{RTB}.

\paragraph{Migliorie per lo \textit{Sprint} 13}
\begin{itemize}
    \item \textbf{Massimizzare lo sforzo:} è necessario il massimo impegno per scongiurare l'accumulo di ulteriore ritardo.
    \item \textbf{Attenta pianificazione:} è necessario rendere i preventivi più accurati. Va incentivata la discussione per ogni compito e vanno introdotte metriche più oggettive per la stima.
    \item \textbf{Comunicazione e allineamento:} ogni mancata comunicazione può essere altro ritardo. Non basta più la garanzia individuale, il \textbf{Responsabile} deve farsi carico di ottenere aggiornamenti frequenti dai membri.
\end{itemize}


\subsubsection{\textit{Sprint} 13: 15/08/2026 - 21/08/2026}

\paragraph{Obiettivi e attività}

Questo \textit{sprint} avrà una lunghezza standard (7 gg). Le attività saranno incentrate sul completamento della progettazione del sistema, prevista da questo documento per mercoledì il 19 agosto, predisponendo contestualmente l'avvio della fase di programmazione. Sarà anche lo \textit{sprint} del colloquio col Prof. Vardanega e di conclusione formale della fase \textbf{RTB}. Di seguito le attività individuate:

\begin{itemize}
    \item Preparazione, richiesta e incontro con il prof. Vardanega, dopo il rinvio dall'ultimo \textit{sprint}.
    \item Termine della progettazione del sistema, in particolare del \textit{backend}, dei \textit{design pattern}, dei test e integrazione delle parti.
    \item Normazione dei processi di implementazione del prodotto: codifica, test, strumenti di \textit{build}, \textit{\glo{workflow}{workflow}}.
\end{itemize}

\paragraph{\textit{Preventivo} orario ed economico}

I ruoli assegnati sono: \textbf{Responsabile}, \textbf{Amministratore}, \textbf{Progettista}, \textbf{Programmatore} e \textbf{Verificatore}. Il focus dello \textit{sprint} è ancora sulla progettazione, ma è previsto e auspicato l'inizio della fase di programmazione.

\begin{table}[H]
    \centering
    \renewcommand{\arraystretch}{1.3}
    \begin{tabularx}{0.9\textwidth}{|X|c|c|c|c|c|c|c|}
        \hline
        \textbf{\textit{Sprint} 13} & \multicolumn{6}{|c|}{\textbf{Ruolo}} & \\ \hline
        \textbf{Membro} & \textbf{RE} & \textbf{AM} & \textbf{AN} & \textbf{PG} & \textbf{PR} & \textbf{VE} & \textbf{Totale} \\ \hline
        Andrea Fistarol & 2 & 1 & & & &  & \textbf{3} \\ \hline
        Dario Pirrone & & & & & & 3 & \textbf{3} \\ \hline
        Edoardo Granziero & & & & 3 & 8 & & \textbf{11} \\ \hline
        Marco Barbiero & & & & & & 7 & \textbf{7} \\ \hline
        Samuele Bertin & & & & 8 & 2 & & \textbf{10} \\ \hline
        Sara Ristovic & & & & 10 & & & \textbf{10} \\ \hline
        \textbf{Ore totali per ruolo} & \textbf{2} & \textbf{1} & \textbf{-} & \textbf{19} & \textbf{10} & \textbf{10} & \textbf{42}\\ \hline
    \end{tabularx}
    \caption{\textit{Preventivo} orario - \textit{sprint} 13}
\end{table}

Le risorse economiche previste per lo \textit{sprint} sono:

\begin{table}[H]
    \centering
    \renewcommand{\arraystretch}{1.3}
    \begin{tabularx}{0.9\textwidth}{|X|c|c|c|c|c|c|c|}
        \hline
        \textbf{\textit{Sprint} 13} & \textbf{RE} & \textbf{AM} & \textbf{AN} & \textbf{PG} & \textbf{PR} & \textbf{VE} & \textbf{Totale} \\ \hline
        \textbf{Costo orario} & \EUR{30} & \EUR{20} & \EUR{25} & \EUR{25} & \EUR{15} & \EUR{15} & / \\ \hline
        \textbf{Costo totale} & \EUR{60} & \EUR{20} & \EUR{0} & \EUR{475} & \EUR{150} & \EUR{150} & \textbf{\EUR{855}} \\ \hline
    \end{tabularx}
    \caption{\textit{Preventivo} economico - \textit{sprint} 13}
\end{table}

\paragraph{\textit{Consuntivo}}~

\begin{table}[H]
    \centering
    \renewcommand{\arraystretch}{1.3}
    \begin{tabularx}{\textwidth}{|c|X|c|}
        \hline
        \textbf{\#} & \textbf{Descrizione} & \textbf{Stato} \\ \hline
        223 & Termine della progettazione del \textit{backend} per soddisfare tutti i requisiti & 80\% \\ \hline
        224 & Progettazione di test esaustivi & 50\% \\ \hline
        225 & Progettazione dell'integrazione tra \textit{backend}, \textit{frontend} e \textit{AWS} & 66\% \\ \hline
        226 & Creazione del grafico \textit{Burndown} per lo \textit{sprint} 13 & Completata \\ \hline
        227 & Scrittura di norme di codifica dettagliate, procedure per i test unitari e di integrazione, scelta e uso degli strumenti di \textit{build} & Completata \\ \hline
        228 & Scrittura norme per le convenzioni \textit{Git Flow} nella gestione dei \textit{branch} e configurazione delle \textit{pipeline} di \textit{CI/CD} & Completata \\ \hline
        229 & Stesura del verbale del 14/08 & Completata \\ \hline
        230 & Stesura del Diario di Bordo del 21/08 & Completata \\ \hline
        231 & Scrittura di \textit{consuntivo} dello \textit{sprint} 12 e \textit{preventivo} dello \textit{sprint} 13 & Completata \\ \hline
        220 & Definizione dei \textit{design pattern} partendo da quelli convalidati nel \textbf{PoC} & Completata \\ \hline
        210 & Preparazione delle slide per l'incontro con il Prof. Vardanega & Completata \\ \hline
        212 & Applicazione delle correzioni del Prof. Cardin all'\textbf{AdR}, portando il documento alla versione 2.0.0 & 0\% \\ \hline
        214 & Applicazione delle modifiche ai documenti secondo le indicazioni del Prof. Vardanega & 15\% \\ \hline
    \end{tabularx}
    \caption{Attività svolte - \textit{sprint 13}}
\end{table}

\subparagraph{Note sui mancati completamenti}
\begin{itemize}
    \item \textbf{223}: non ultimato per sottostima del lavoro necessario.
    \item \textbf{224}: fermato per mancanza di dettagli progettuali.
    \item \textbf{225}: in attesa del completamento di \textbf{223}.
    \item \textbf{212}: non iniziato per bassa priorità.
    \item \textbf{214}: corretto solo il template dei verbali, resto non fatto per bassa priorità.
\end{itemize}

\begin{table}[H]
    \hspace*{-0.7cm}
    \renewcommand{\arraystretch}{1.3}
    \begin{tabular}{|l
            |>{\centering\arraybackslash}p{0.1\textwidth}
            |>{\centering\arraybackslash}p{0.1\textwidth}
            |>{\centering\arraybackslash}p{0.1\textwidth}
            |>{\centering\arraybackslash}p{0.1\textwidth}
            |>{\centering\arraybackslash}p{0.1\textwidth}
            |>{\centering\arraybackslash}p{0.1\textwidth}
            |c c|}
        \hline
        \textbf{\textit{Sprint} 13} & \multicolumn{6}{|c|}{\textbf{Ruolo}} & \multicolumn{2}{|c|}{} \\ \hline
        \textbf{Membro} & \textbf{RE} & \textbf{AM} & \textbf{AN} & \textbf{PG} & \textbf{PR} & \textbf{VE} & \multicolumn{2}{|c|}{\textbf{Totale}} \\ \hline
        A. Fistarol & 4 (+2) & 1 & - & - & - & - & \textbf{5} & (+2) \\ \hline
        D. Pirrone & - & - & - & - & - & 2 (-1) & \textbf{2} & (-1) \\ \hline
        E. Granziero & - & - & - & 3 & 10 (+2) & - & \textbf{13} & (+2) \\ \hline
        M. Barbiero & - & - & - & - & - & 7 & \textbf{7} & (0) \\ \hline
        S. Bertin & - & - & - & 10 (+2) & 0 (-2) & - & \textbf{10} & (0) \\ \hline
        S. Ristovic & - & - & - & 10 (+2) & - & - & \textbf{10} & (+2) \\ \hline
        \textbf{Ore totali} & \textbf{4 (+2)} & \textbf{1} & \textbf{0} & \textbf{23 (+4)} & \textbf{10} & \textbf{9 (-1)} & \textbf{47} & (+5) \\ \hline
    \end{tabular}
    \caption{\textit{Consuntivo} orario - \textit{sprint} 13}
\end{table}

\begin{table}[H]
    \centerline{
        \centering
        \renewcommand{\arraystretch}{1.3}
        \begin{tabularx}{1.00\textwidth}{|X|c|c|c|c|c|c|c|}
            \hline
            \textbf{\textit{Sprint} 13} & \textbf{RE} & \textbf{AM} & \textbf{AN} & \textbf{PG} & \textbf{PR} & \textbf{VE} & \textbf{Totale} \\ \hline  
            \textbf{Costo orario} & \EUR{30} & \EUR{20} & \EUR{25} & \EUR{25} & \EUR{15} & \EUR{15} & / \\ \hline
            \textbf{Costo totale} & \EUR{120} & \EUR{20} & \EUR{0} & \EUR{575} & \EUR{150} & \EUR{135} & \textbf{\EUR{1000}} \\ \hline
            \textbf{{\textDelta} preventivo} & \EUR{+60} & \EUR{0} & \EUR{0} & \EUR{+100} & \EUR{0} & \EUR{-15} & \textbf{\EUR{+145}} \\ \hline
            \textbf{Saldo a fine \textit{sprint}} & \textbf{\EUR{240}} & \textbf{\EUR{80}} & \textbf{\EUR{-150}} & \textbf{\EUR{650}} & \textbf{\EUR{1.245}} & \textbf{\EUR{885}} & \textbf{\EUR{2.950}}\\ \hline
        \end{tabularx}
    }
    \caption{\textit{Consuntivo} economico - \textit{sprint} 13}
\end{table}

\paragraph{Retrospettiva}~\\
Il grado di completamento delle attività continua ad essere aleatorio. A parte i test, che comunque non sono ultimati, i progettisti non sono riusciti a concludere il loro lavoro entro il 19 e iniziare a porre le basi per la codifica. La ragione è sempre l'incapacità di stimare la durata dei compiti, dovuta alla mancata discussione su come questi debbano essere svolti. Il ciclo è sempre lo stesso: si decide cosa fare, lo si assegna e poi ognuno cerca di capire come deve svolgere il proprio compito, prima e durante l'esecuzione. Tuttavia, in fase di pianificazione, ogni stima è completamente inutile, non si ha idea di cosa ci sia da fare nella pratica. Un ulteriore problema causato da questo modo di procedere è il rischio di ritrovarsi con prodotti non idonei in fase di verifica o ancora successivamente, quando ci si può rendere conto che quanto fatto da un membro non sia adeguato o ai fini o alla prassi, con annessi i costi che una correzione può comportare. Alla luce di un ritardo già molto significativo, se ne sta accumulando altro.

Necessaria la normalizzazione del saldo economico, cioè della somma negativa del ruolo \textbf{Analista}, in quanto il ruolo e lo sforamento delle ore previste dal \textit{preventivo} iniziale non son più rilevanti in questa fase del progetto.

\begin{figure}[H]
  \centering
    \includegraphics[width=0.9\textwidth]{grafici/burndown_sprint13.png}
    \caption{\textit{Grafico Burndown} - \textit{sprint} 13}
    \label{fig:burndown-sprint13}
\end{figure}

\paragraph{Misure Correttive}
\begin{itemize}
    \item Analizzare e programmare le restanti attività di progettazione per concluderle entro il 23/08.
    \item Ogni membro deve arrivare alla riunione di fine \textit{sprint} (il 31/08) con già idea delle prossime attività.
    \item Dedicare più tempo alla discussione delle attività e assicurarsi che oltre al cosa sia definito il come, anteriormente all'esecuzione.
\end{itemize}

\subparagraph{Analisi dei Rischi}
\begin{itemize}
    \item \textbf{RO1, RO4}: ricorrenti i rischi di Scarsa Pianificazione e Sottostima delle Attività. Le precauzioni si dimostrano inefficaci o mal applicate, la loro pericolosità invece aumenta per il ritardo accumulato - quella di RO4 è stata aumentata a Alta.
    \item \textbf{(Nuovo) RO8 - Riunioni Inefficaci}: un altro aspetto dei rischi sopra è che le riunioni preposte alla pianificazione, quelle di fine/inizio \textit{sprint}, non sono efficaci. La loro durata non è proporzionata alla necessità ma sottoposta alla disponibilità dei membri. Se le riunioni vengono viste come un peso, allora diventa un peso anche una corretta organizzazione, e ciò non è possibile. È essenziale trarre il massimo giovamento da ogni occasione di discussione interna e quelle di pianificazione devono essere efficaci, pena il ripresentarsi di rischi ricorrenti. D'ora in poi andrà concordato in anticipo un Ordine del Giorno (OdG) per ogni riunione e ogni punto andrà affrontato. Ciascun membro dovrà prepararsi in modo appropriato alla discussione e dovrà dedicarvi il suo tempo, rispetto a una durata della riunione congrua al numero e al peso dei punti all'OdG.
\end{itemize}

\paragraph{Miglioramento della pianificazione futura}
\subparagraph{Breve termine}
\begin{itemize}
    \item Da aspettarsi un ritardo di 4 gg per la conclusione della fase di codifica, precedentemente prevista al 31/08, ora ragionevole aspettarsi al 4/09. Prima di agire, monitorare lo sviluppo durante lo \textit{sprint}.
\end{itemize}

\subparagraph{Lungo termine}
\begin{itemize}
    \item Riallineamento del \textit{preventivo} a finire saldando il debito dell'\textbf{Analista} e spostando ore in eccesso dal \textbf{Progettista} agli altri ruoli in carenza, come \textbf{Responsabile} e \textbf{Amministratore}.
    \item La data della \textit{milestone} "Test e raffinamento" è da posticipare di 7 gg all'11/09, più realistica dati i ritardi.
    \item La data di consegna (21/09), già pessimistica, può essere mantenuta.
\end{itemize}

\paragraph{Preventivo a Finire}~
\begin{table}[H]
    \centering
    \renewcommand{\arraystretch}{1.3}
    \begin{tabularx}{\textwidth}{|X|c|c|c|c|c|c|c|}
        \hline
         \textbf{Per ruolo}& \textbf{RE} & \textbf{AM} & \textbf{AN} & \textbf{PG} & \textbf{PR} & \textbf{VE} & \textbf{Totale} \\ \hline
         \textbf{Ore} & 11 & 13 & 0 & 3 & 86 & 66 & 179 \\ \hline
         \textbf{Costo} & \EUR{330} & \EUR{260} & \EUR{0} & \EUR{75} & \EUR{1.290} & \EUR{990} & \EUR{2.945} \\ \hline \hline
         \textbf{{\textDelta} PaF prec.} & \textbf{RE} & \textbf{AM} & \textbf{AN} & \textbf{PG} & \textbf{PR} & \textbf{VE} & \textbf{Totale} \\ \hline
         \textbf{{\textDelta} ore} & +3 & +9 & +6 & -23 & +3 & +7 & +5 \\ \hline
         \textbf{{\textDelta} costo} & \EUR{+90} & \EUR{+180} & \EUR{+150} & \EUR{-575} & \EUR{+45} & \EUR{+105} & \EUR{-5} \\ \hline
    \end{tabularx}
    \caption{Nuovo \textit{Preventivo} a Finire - \textit{sprint} 13}
\end{table}

\begin{table}[H]
    \centering
    \renewcommand{\arraystretch}{1.3}
    \begin{tabularx}{0.7\textwidth}{|X|c|}
        \hline
         \textbf{Membro} & \textbf{Ore rimaste (su 94)} \\ \hline
         Andrea Fistarol & 45 \\ \hline
         Dario Pirrone & 20 \\ \hline
         Edoardo Granziero & 38 \\ \hline
         Marco Barbiero & 21 \\ \hline
         Samuele Bertin & 25 \\ \hline
         Sara Ristovic & 25 \\ \hline
    \end{tabularx}
    \caption{Ore rimaste - \textit{sprint} 13}
\end{table}

\subsection{\textit{Sprint} 14: 22/08/2026 - 31/08/2026}
\paragraph{Obbiettivi e attività} Lo \textit{sprint} durerà 10 gg e il suo termine corrisponde con quello previsto per la fase di codifica dell'\textbf{MVP}.
Infatti, il focus delle attività è appunto la traduzione in codice delle scelte progettuali. Seguono le attività programmate.
\begin{itemize}
    \item Conclusione della progettazione del \textit{backend} e integrazione progettuale, entro e non oltre domenica 23.
    \item Codifica di \textit{frontend}, \textit{backend}, \textit{DB} e agenti.
    \item Configurazione dei servizi \textit{AWS}.
    \item Predisposizione dei \textit{workflow} di \textit{CI/CD} e della \textit{repository} per l'\textbf{MVP}.
\end{itemize}

\paragraph{\textit{Preventivo} orario ed economico}

\begin{table}[H]
    \centering
    \renewcommand{\arraystretch}{1.3}
    \begin{tabularx}{0.9\textwidth}{|X|c|c|c|c|c|c|c|}
        \hline
        \textbf{\textit{Sprint} 14} & \multicolumn{6}{|c|}{\textbf{Ruolo}} & \\ \hline
        \textbf{Membro} & \textbf{RE} & \textbf{AM} & \textbf{AN} & \textbf{PG} & \textbf{PR} & \textbf{VE} & \textbf{Totale} \\ \hline
        Andrea Fistarol & 5 & 5 & - & - & - & -  & \textbf{10} \\ \hline
        Dario Pirrone & - & - & - & - & 8 & - & \textbf{8} \\ \hline
        Edoardo Granziero & - & - & - & - & - & 8 & \textbf{8} \\ \hline
        Marco Barbiero & - & - & - & - & 15 & - & \textbf{15} \\ \hline
        Samuele Bertin & - & - & - & - & 20 & - & \textbf{20} \\ \hline
        Sara Ristovic & - & - & - & 2 & 11 & - & \textbf{13} \\ \hline
        \textbf{Ore totali per ruolo} & \textbf{5} & \textbf{5} & \textbf{0} & \textbf{2} & \textbf{54} & \textbf{8} & \textbf{74}\\ \hline
    \end{tabularx}
    \caption{\textit{Preventivo} orario - \textit{sprint} 14}
\end{table}

Le risorse economiche previste per lo \textit{sprint} sono:

\begin{table}[H]
    \centering
    \renewcommand{\arraystretch}{1.3}
    \begin{tabularx}{0.9\textwidth}{|X|c|c|c|c|c|c|c|}
        \hline
        \textbf{\textit{Sprint} 14} & \textbf{RE} & \textbf{AM} & \textbf{AN} & \textbf{PG} & \textbf{PR} & \textbf{VE} & \textbf{Totale} \\ \hline
        \textbf{Costo orario} & \EUR{30} & \EUR{20} & \EUR{25} & \EUR{25} & \EUR{15} & \EUR{15} & / \\ \hline
        \textbf{Costo totale} & \EUR{150} & \EUR{100} & \EUR{0} & \EUR{50} & \EUR{810} & \EUR{120} & \textbf{\EUR{1.230}} \\ \hline
    \end{tabularx}
    \caption{\textit{Preventivo} economico - \textit{sprint} 14}
\end{table}

\paragraph{\textit{Consuntivo}}~

\begin{table}[H]
    \centering
    \renewcommand{\arraystretch}{1.3}
    \begin{tabularx}{\textwidth}{|c|X|c|}
        \hline
        \textbf{\#} & \textbf{Descrizione} & \textbf{Stato} \\ \hline
        223 & Termine della progettazione del \textit{backend} per soddisfare tutti i requisiti & Completata \\ \hline
        224 & Progettazione di test esaustivi & 90\% \\ \hline
        225 & Progettazione dell'integrazione tra \textit{backend}, \textit{frontend} e \textit{AWS} & Completata \\ \hline
        237 & Creazione del grafico \textit{Burndown} per lo \textit{sprint} 14 & 0\% \\ \hline
        236 & Creazione slide per incontro di monitoraggio & Completata \\ \hline
        238 & Stesura del verbale del 21/08 & Completata \\ \hline
        239 & Scrittura di \textit{consuntivo} dello \textit{sprint} 13 e \textit{preventivo} dello \textit{sprint} 14 & Completata \\ \hline
        -- & \textit{Issue} riguardanti alla programmazione dell'infrastruttura \textit{AWS} & Completata \\ \hline
        -- & \textit{Issue} riguardanti alla programmazione del \textit{frontend} & Completata \\ \hline
        -- & \textit{Issue} riguardanti alla programmazione del \textit{backend} & 85\% \\ \hline
        -- & \textit{Issue} riguardanti alla programmazione degli agenti & 95\% \\ \hline
        278 & Glossario: Aggiungere build da \textbf{PdP} e collegarlo & 0\% \\ \hline
        279 & Glossario: Inserire i nuovi termini delle \textbf{NdP} e collegarli & 0\% \\ \hline
        280 & Re-baselining del preventivo orario: riallineamento risorse e ruoli & Completata \\ \hline
        281 & \textbf{NdP}: Aggiornare struttura repo & Completata \\ \hline
        283 & Setup dei workflow (\textit{PR} check, deploy staging, release production) previa revisione critica & Completata \\ \hline
        284 & Predisposizione della configurazione dei test con pnpm (Vitest, pytest) & Completata \\ \hline
        287 & Aggiornamento repository secondo \textbf{NdP} (\textit{branch} e protezione) & Completata \\ \hline
        285 & Predisposizione della configurazione del linting e del formatting con pnpm (Biome, Pylint, Black) & Completata \\ \hline
        212 & Applicazione delle correzioni del Prof. Cardin all'\textbf{AdR}, portando il documento alla versione 2.0.0 & 0\% \\ \hline
        214 & Applicazione delle modifiche ai documenti secondo le indicazioni del Prof. Vardanega & 60\% \\ \hline
    \end{tabularx}
    \caption{Attività svolte - \textit{sprint 14}}
\end{table}

\subparagraph{Note sui mancati completamenti}
\begin{itemize}
    \item \textbf{224}: valutazione errata del soddifacimento della \textit{issue}.
    \item \textbf{237}: a causa del mancato tracciamento corretto del andamento delle \textit{issue} non è stato possibile realizzare un grafico utile.
    \item \textbf{programmazione backend}: sottostima del lavoro.
    \item \textbf{programmazione agenti}: sottostima del lavoro e conflitto con \textit{issue} di progettazione test (le issue mancanti riguardano i test degli agenti).
    \item \textbf{278}: sottostima del lavoro.
    \item \textbf{279}: sottostima del lavoro.
    \item \textbf{212}: non iniziata per bassa priorità.
    \item \textbf{214}: applicate alcune modifiche ai documenti, resto non fatto per bassa priorità, è necessaio dividere il restante lavoro in \textit{issue} più atomiche.
\end{itemize}

\begin{table}[H]
    \hspace*{-0.7cm}
    \renewcommand{\arraystretch}{1.3}
    \begin{tabular}{|l
            |>{\centering\arraybackslash}p{0.1\textwidth}
            |>{\centering\arraybackslash}p{0.1\textwidth}
            |>{\centering\arraybackslash}p{0.1\textwidth}
            |>{\centering\arraybackslash}p{0.1\textwidth}
            |>{\centering\arraybackslash}p{0.1\textwidth}
            |>{\centering\arraybackslash}p{0.1\textwidth}
            |c c|}
        \hline
        \textbf{\textit{Sprint} 14} & \multicolumn{6}{|c|}{\textbf{Ruolo}} & \multicolumn{2}{|c|}{} \\ \hline
        \textbf{Membro} & \textbf{RE} & \textbf{AM} & \textbf{AN} & \textbf{PG} & \textbf{PR} & \textbf{VE} & \multicolumn{2}{|c|}{\textbf{Totale}} \\ \hline
        A. Fistarol & 4 (-1) & 3 (-2) & - & - & - & 1 (+1) & \textbf{8} & (-2) \\ \hline
        D. Pirrone & - & - & - & - & 8 & - & \textbf{8} & (0) \\ \hline
        E. Granziero & - & - & - & - & - & 6 (-2) & \textbf{6} & (-2) \\ \hline
        M. Barbiero & - & - & - & - & 14 (-1) & - & \textbf{14} & (-1) \\ \hline
        S. Bertin & - & - & - & - & 14 (-6) & - & \textbf{14} & (-6) \\ \hline
        S. Ristovic & - & - & - & 1 (-1) & 12 (+1) & - & \textbf{13} & (0) \\ \hline
        \textbf{Ore totali} & \textbf{4 (-1)} & \textbf{3 (-2)} & \textbf{0} & \textbf{1 (-1)} & \textbf{48 (-6)} & \textbf{7 (-1)} & \textbf{63} & (+5) \\ \hline
    \end{tabular}
    \caption{\textit{Consuntivo} orario - \textit{sprint} 14}
\end{table}

\begin{table}[H]
    \centerline{
        \centering
        \renewcommand{\arraystretch}{1.3}
        \begin{tabularx}{1.00\textwidth}{|X|c|c|c|c|c|c|c|}
            \hline
            \textbf{\textit{Sprint} 14} & \textbf{RE} & \textbf{AM} & \textbf{AN} & \textbf{PG} & \textbf{PR} & \textbf{VE} & \textbf{Totale} \\ \hline  
            \textbf{Costo orario} & \EUR{30} & \EUR{20} & \EUR{25} & \EUR{25} & \EUR{15} & \EUR{15} & / \\ \hline
            \textbf{Costo totale} & \EUR{120} & \EUR{60} & \EUR{0} & \EUR{25} & \EUR{720} & \EUR{105} & \textbf{\EUR{1030}} \\ \hline
            \textbf{{\textDelta} preventivo} & \EUR{-30} & \EUR{-40} & \EUR{0} & \EUR{-25} & \EUR{-90} & \EUR{-15} & \textbf{\EUR{-200}} \\ \hline
            \textbf{Saldo a fine \textit{sprint}} & \textbf{\EUR{210}} & \textbf{\EUR{200}} & \textbf{\EUR{0}} & \textbf{\EUR{50}} & \textbf{\EUR{570}} & \textbf{\EUR{885}} & \textbf{\EUR{1.915}}\\ \hline
        \end{tabularx}
    }
    \caption{\textit{Consuntivo} economico - \textit{sprint} 14}
\end{table}

\paragraph{Retrospettiva}~\\
Al termine dello \textit{sprint} non tutte le \textit{issue} sono state portate a termine, questo a causa di una mancata analisi dei tempi necessari per svolgere ciascuna \textit{issue}; in questo \textit{sprint} si è cercato di avere una pianificazione di \textit{issue} più atomiche e con una stima preventiva dei tempi per ciascuna \textit{issue} oltre a delle scadenze specifiche, questo approccio sicuramente è un passo verso la giusta direzione ma senza una corretta istanziazione di tempi di analisi del lavoro effettivo da svolgere prima della creazione ed assegnazione delle \textit{issue} il problema di stima errata dei tempi di lavoro rimane presente, inoltre il modo in cui sono state strutturate le \textit{issue} in questo \textit{sprint} ha portato un contrasto con le nuove norme di progetto instanziate ad inizio \textit{sprint}, in particolare le norme relative ai \textit{branch} e ai \textit{commit} su \textit{GitHub}.



\paragraph{Misure Correttive}
\begin{itemize}
    \item Ogni membro del gruppo è chiamato a fermare i lavori in sviluppo almeno 24 ore prima della riunione di fine \textit{sprint}, per dedicarsi totalmente nella analisi del lavoro restante e pianificazione di nuove \textit{issue}, in particolare riguardanti gli aspetti di progetto di cui si ha maggiore conoscenza.
\end{itemize}

\subparagraph{Analisi dei Rischi}
\begin{itemize}
    \item \textbf{RO1, RO4}: Restano sempre presenti i rischi di scarsa organizzazione e sottostima delle attività, questi rischi risultano complessi da mitigare anche a causa del tempo estremamente ristretto prima della conclusione del progetto, in queste condizioni è complesso trovare tempi extra per pianificare i compiti con la dovuta cura.
    \item \textbf{RO8 - Riunioni Inefficaci}: Le riunioni interne hanno durata superiore al necessario portando a stanchezza e perdita di focus nelle decisione prese.
\end{itemize}



\subsection{\textit{Sprint} 15: 01/09/2026 - 7/09/2026}
\paragraph{Obbiettivi e attività} Lo \textit{sprint} durerà 7 gg e il suo termine sarà un giorno prima del incontro di esposizione del \textbf{MVP} al azienda \textbf{proponente} (8 settembre, data di "backup" il 10 settembre).
Il focus delle attività pianificate riguarda il completamento del \textbf{MVP}, con i dovuti test, più nel dettaglio:
\begin{itemize}
    \item Conclusione della programmazione del \textit{backend}.
    \item Allineamento punti di contatto tra \textit{backend}, \textit{frontend} ed infrastruttura \textit{AWS}.
    \item Debug e test del \textbf{MVP}.
    \item Stesura \textbf{Specifiche Tecniche}.
    \item Stesura \textbf{Manuale Utente}.
\end{itemize}

\paragraph{\textit{Preventivo} orario ed economico}

\begin{table}[H]
    \centering
    \renewcommand{\arraystretch}{1.3}
    \begin{tabularx}{0.9\textwidth}{|X|c|c|c|c|c|c|c|}
        \hline
        \textbf{\textit{Sprint} 15} & \multicolumn{6}{|c|}{\textbf{Ruolo}} & \\ \hline
        \textbf{Membro} & \textbf{RE} & \textbf{AM} & \textbf{AN} & \textbf{PG} & \textbf{PR} & \textbf{VE} & \textbf{Totale} \\ \hline
        Andrea Fistarol & - & - & - & - & 17 & -  & \textbf{17} \\ \hline
        Dario Pirrone & 2 & - & - & - & - & 5 & \textbf{7} \\ \hline
        Edoardo Granziero & - & - & - & - & 15 & - & \textbf{15} \\ \hline
        Marco Barbiero & - & - & - & - & - & - & \textbf{0} \\ \hline
        Samuele Bertin & - & 2 & - & - & 3 & - & \textbf{5} \\ \hline
        Sara Ristovic & - & - & - & - & 8 & - & \textbf{8} \\ \hline
        \textbf{Ore totali per ruolo} & \textbf{2} & \textbf{2} & \textbf{0} & \textbf{0} & \textbf{43} & \textbf{5} & \textbf{52}\\ \hline
    \end{tabularx}
    \caption{\textit{Preventivo} orario - \textit{sprint} 15}
\end{table}

Le risorse economiche previste per lo \textit{sprint} sono:

\begin{table}[H]
    \centering
    \renewcommand{\arraystretch}{1.3}
    \begin{tabularx}{0.9\textwidth}{|X|c|c|c|c|c|c|c|}
        \hline
        \textbf{\textit{Sprint} 15} & \textbf{RE} & \textbf{AM} & \textbf{AN} & \textbf{PG} & \textbf{PR} & \textbf{VE} & \textbf{Totale} \\ \hline
        \textbf{Costo orario} & \EUR{30} & \EUR{20} & \EUR{25} & \EUR{25} & \EUR{15} & \EUR{15} & / \\ \hline
        \textbf{Costo totale} & \EUR{60} & \EUR{40} & \EUR{0} & \EUR{0} & \EUR{645} & \EUR{75} & \textbf{\EUR{820}} \\ \hline
    \end{tabularx}
    \caption{\textit{Preventivo} economico - \textit{sprint} 15}
\end{table}
\end{document}